% !TEX root = ../main.tex

\begin{story}

以下この節において定理3.2の証明を行う。ただし、細かい計算には立ち入らず、筋道を示すことを第1の目標としよう。記号等は前節のものを踏襲して用いる。

まず，ハミルトニアン $H_l$ の形(3.73)を，フロベニウスの方法により決定する
\end{story}

\begin{Q}[【核心】「ハミルトニアン $H_l$ の形を、フロベニウスの方法により決定する」とは具体的にどういうことか？]
\textbf{A.} これは、前節までの「見かけの特異点 $x=\lambda_k$（特性指数 $0, 2$）」における局所的な解の振る舞いから、系全体を統制するエネルギー関数（ハミルトニアン）を逆算して導き出す、という驚くべきプロセスを指しています。
「フロベニウスの方法」とは、特異点周りで微分方程式の解を冪級数展開して求める標準的な手法です。特性指数の差が整数の場合、通常は解に「対数項（$\log(x-\lambda_k)$）」が現れます。しかし、$x=\lambda_k$ は「見かけの特異点」であるため、この対数項の係数は厳密に $0$ にならなければなりません。
この「対数項の係数が $0$ になる」という代数的な条件式を書き下すと、そこに変数 $\lambda_k$ と留数 $\mu_k$、そして変形パラメータ $t_l$ が結びついた関係式が浮かび上がります。この関係式を整理して得られる多項式（または有理関数）こそが、ハミルトニアン $H_l$ そのものになるのです。局所的な「対数項回避の条件」が、大域的な「ハミルトン力学系の生成子」を決定するというのが、岡本理論の最も美しいハイライトの一つです。
\end{Q}

\begin{Q}[【深掘り】そもそも「ガルニエ系 (Garnier system)」とは何か？]
\textbf{A.} 一言で言えば、「パンルヴェ方程式の多変数化（高次元化）」です。
パンルヴェ方程式は、独立変数が $t$ の1つだけ（方程式の特異点が例えば $0, 1, t, \infty$ の4つ）である場合のモノドロミー保存変形を記述する非線形常微分方程式でした。
ガルニエ系は、動かせる特異点（変形パラメータ）を $t_1, t_2, \dots, t_g$ と複数個に増やしたときに現れる、非線形偏微分方程式系です。変数が多次元になっても、見かけの特異点 $(\lambda_k, \mu_k)$ を正準変数とみなすことで、時間発展 $\partial/\partial t_l$ がハミルトニアン $H_l$ によって完全に統制される「多変数ハミルトン系」として綺麗に書き表すことができます。この節では、そのハミルトニアンが実際にどのような数式になるかを導出していくことになります。
\end{Q}

\begin{story}
$x=\lambda_k$ を微分方程式 (3.33) の見かけの特異点とし、局所座標 $u = x - \lambda_k$ をとる。$u=0$ の近傍で (3.33) を
\begin{equation}
    u^2 \frac{d^2 y}{du^2} + q_1(u,t)u \frac{dy}{du} + q_2(u,t)y = 0 \label{eq:3.77}
\end{equation}
と書く。確定特異点 $u=0$ における特性指数は $0, 2$ であり、これが非対数的であるための条件は、ベキ級数解
\begin{equation}
    y = 1 + y_1 u + y_2 u^2 + \cdots \label{eq:3.78}
\end{equation}
が存在することである。

\begin{proposition}[対数項を持たないための適合条件]
\eqref{eq:3.77} の係数を次のように展開する。
\begin{equation*}
    q_1(u,t) = -1 + a_k u + a'_k u^2 + \cdots, \quad q_2(u,t) = \mu_k u + b_k u^2 + b'_k u^3 + \cdots
\end{equation*}
この表示と \eqref{eq:3.78} を \eqref{eq:3.77} に代入し、未定係数法により比較すると、
\begin{equation*}
    y_1 = \mu_k, \quad j(j-2)y_j = R_j(y_1, \cdots, y_{j-1}) \quad (j=2, \cdots)
\end{equation*}
となるから、ベキ級数解 \eqref{eq:3.78} が存在するための条件は、$j=2$ の場合を考えて
\begin{equation}
    R_2(y_1) = \mu_k y_1 + a_k y_1 + b_k = \mu_k^2 + a_k \mu_k + b_k = 0 \label{eq:3.79}
\end{equation}
であることがわかる。
\end{proposition}

\end{story}

\begin{Q}[【核心】なぜ $j(j-2)y_j = R_j$ において $j=2$ のとき、$R_2(y_1) = 0$ が「対数項を持たない条件」になるのか？]
\textbf{A.} これはフロベニウス法の決定方程式の構造に由来します。特性指数が $\rho = 0, 2$ ということは、微分方程式の左辺に $u^\rho$ を代入したときの最低次数の係数（決定方程式）が $\rho(\rho-2)=0$ になるということです。
特性指数 $\rho=0$ に対応する級数解を求めようとして係数 $y_j$ を順に決めていくと、漸化式の左辺は $j(j-2)y_j$ となります。$j=1$ のときは $1 \cdot (-1) y_1 = R_1$ なので $y_1$ は一意に決まりますが、$j=2$ のステップに到達した瞬間、左辺は $2(0)y_2 = 0 \cdot y_2 = 0$ となってしまいます。
もしこのとき右辺の $R_2$ が $0$ でなければ、「$0 = (\text{0でない値})$」という矛盾が生じ、$y_2$ を決定することができません。この矛盾を回避して解を構成するためには、どうしても独立な解として $u^2 \log u$ のような「対数項」を付け加える必要が出てきます。
逆に言えば、対数項が現れない（見かけの特異点である）ためには、どうしても右辺が $R_2 = 0$ になってくれて、$0 \cdot y_2 = 0$（$y_2$ は任意定数としてよい）となる必要があるのです。この $R_2=0$ こそが、特異点が見かけ上のものであるための強烈な代数的制約（適合条件）です。
\end{Q}

\begin{story}
\begin{theorem}[ハミルトニアンを決定する連立1次方程式]
一方、元の方程式の係数を用いて $a_k, b_k$ を具体的に計算すると、
\begin{align*}
    a_k &= \frac{1-\kappa_0}{\lambda_k} + \frac{1-\kappa_1}{\lambda_k-1} + \sum_{l=1}^g \frac{1-\theta_l}{\lambda_k-t_l} - \sum_{p=1(p \neq k)}^g \frac{1}{\lambda_k-\lambda_p} \\
    b_k &= \frac{\kappa}{\lambda_k(\lambda_k-1)} + \sum_{p=1(p \neq k)}^g \frac{\lambda_p(\lambda_p-1)\mu_p}{\lambda_k(\lambda_k-1)(\lambda_k-\lambda_p)} - \frac{2\lambda_k-1}{\lambda_k(\lambda_k-1)}\mu_k - \sum_{l=1}^g \frac{t_l(t_l-1)H_l}{\lambda_k(\lambda_k-1)(\lambda_k-t_l)}
\end{align*}
である。これを \eqref{eq:3.79} に代入して整理した式は、
\begin{equation}
    \sum_{l=1}^g E_{kl} H_l = \mu_k^2 + \hat{a}_k \mu_k + \hat{b}_k \label{eq:3.80}
\end{equation}
\begin{equation*}
    E_{kl} = \frac{t_l(t_l-1)}{\lambda_k(\lambda_k-1)(\lambda_k-t_l)}
\end{equation*}
という形をしているから、これを $H_l$ に関する連立1次方程式と思って解けばよい。
\end{theorem}
\end{story}

\begin{Q}[【深掘り】この連立方程式 \eqref{eq:3.80} が解けることで、理論上どのような「嬉しいこと」があるのか？]
\textbf{A.} これは「ガルニエ系のハミルトニアン $H_l$ が、正準変数 $(\lambda_k, \mu_k)$ の有理関数として具体的に構成できる」という決定的な事実を示しています。
左辺の係数行列 $E_{kl}$ はコーシー行列の一種であり、$\lambda_k$ と $t_l$ がすべて相異なるという仮定のもとで正則（逆行列を持つ）になります。したがって、この連立1次方程式は必ず一意な解 $H_l$ を持ちます。
局所的な「対数項を消去せよ」という解析的な条件 $R_2=0$ を解きほぐしただけで、系全体の時間発展 $\partial/\partial t_l$ を支配する力学系のエネルギー関数 $H_l$ が代数的に完全に決定されてしまうという、モノドロミー保存変形の理論構造の美しさがここに凝縮されています。
\end{Q}

\begin{story}
\begin{lemma}[3.4]
行列 $E = ((E_{kl}))$ の逆行列を $F = ((F_{lk}))$ とすると、
\begin{equation}
    F_{lk} = M_l M^{k,l} \label{eq:3.81}
\end{equation}
\begin{equation*}
    M_l = - \frac{L(t_l)}{T'(t_l)}, \quad M^{k,l} = \frac{T(\lambda_k)}{L'(\lambda_k)(\lambda_k - t_l)}
\end{equation*}
となる。
\end{lemma}

補題に出てくる、$t, \lambda, \mu$ の有理関数 $M_l, M^{k,l}$ は、それぞれ (3.74), (3.75) で与えられたものである。(3.81) により、(3.80) を $H_l$ について解けば (3.73) が得られる。そのとき、関係式
\begin{equation*}
    \sum_{k=1}^g \frac{M^{k,l}}{\lambda_k(\lambda_k-1)} = 1
\end{equation*}
などを用いて計算することになるが、(3.73) を導く計算の詳細は省略する。

\begin{proof}
線型代数の演習問題であるから、簡単に証明しておこう。
$x$ の有理関数
\begin{equation}
    Z_l(x) = \frac{T(x)}{x(x-1)(x-t_l)L(x)}
\end{equation}
を考える。ここで、$T(x), L(x)$ は、(3.69), (3.70) で与えられる $x$ の多項式、すなわち
\begin{equation*}
    T(x) = x(x-1)\prod_{m=1}^g (x-t_m), \quad L(x) = \prod_{k=1}^g (x-\lambda_k)
\end{equation*}
である。$Z_l(x)$ は $x=\lambda_k \ (k=1, \cdots, g)$ に、1位の極をもつ。さらに、有理関数 $Z_l(x)$ の分母は $g$ 次多項式、分子は $g-1$ 次多項式、であることに注意してこの有理関数を部分分数展開する。実際
\begin{equation*}
    \underset{x=\lambda_k}{\mathrm{Res}} Z_l(x) = \frac{T(\lambda_k)}{\lambda_k(\lambda_k-1)(\lambda_k-t_l)L'(\lambda_k)} = \frac{M^{k,l}}{\lambda_k(\lambda_k-1)}
\end{equation*}
より、次の式が得られる。
\begin{equation}
    Z_l(x) = \sum_{k=1}^g \frac{M^{k,l}}{\lambda_k(\lambda_k-1)} \cdot \frac{1}{x-\lambda_k}
\end{equation}
ここで、まず
\begin{equation*}
    Z_l(t_h) = \sum_{k=1}^g \frac{M^{k,l}}{\lambda_k(\lambda_k-1)} \cdot \frac{1}{t_h-\lambda_k} = - \sum_{k=1}^g \frac{M^{k,l} E_{kh}}{t_h(t_h-1)}
\end{equation*}
となることに気がつく．他方，定義式から
\begin{equation*}
    Z_l(t_l) = \frac{T'(t_l)}{t_l(t_l-1)L(t_l)} = \frac{-1}{t_l(t_l-1)M_l}
\end{equation*}
\begin{equation*}
    Z_l(t_h) = 0 \quad (h \neq l)
\end{equation*}
であるから，結局
\begin{align*}
    1 &= \sum_{k=1}^g E_{kl} M_l M^{k,l} \\
    0 &= \sum_{k=1}^g E_{kh} M_l M^{k,l} \quad (h \neq l)
\end{align*}
が得られる．これは(3.81)を意味している．

(3.80), (3.81)から $H_l$ を求めれば(3.73)が得られる。
\end{proof}
\end{story}

\begin{Q}[【核心】なぜ逆行列を求めるために、突然 $Z_l(x)$ という有理関数を考えるのか？]
\textbf{A.} 行列 $E$ と $F$ が互いに逆行列であること（すなわち $\sum E_{kl}F_{lm} = \delta_{km}$）を直接計算で証明するのは、要素が複雑な分数であるため非常に困難です。
しかし、$Z_l(x)$ の定義式に $T(x)$ を代入して約分すると、実は
$$ Z_l(x) = \frac{\prod_{m \neq l}^g (x-t_m)}{L(x)} $$
となり、これは分子が $g-1$ 次、分母が $g$ 次の有理関数になります。
複素解析において「分子の次数が分母の次数より2以上小さい有理関数の、すべての極における留数の和はゼロになる」という強力な定理（留数定理）があります。この定理や、部分分数展開に $x=t_m$ などを代入して得られる「ラグランジュ補間の恒等式」を利用することで、行列の積の和 $\sum$ を「留数の和」にすり替えることができ、複雑な代数計算を一瞬でクロネッカーのデルタ $\delta_{km}$ に帰着させることができるのです。
\end{Q}

\begin{Q}[【深掘り】多項式 $T(x)$ と $L(x)$ はそれぞれ何を意味しているのか？]
\textbf{A.} 多変数（$g$ 変数）のガルニエ系において、多数の特異点を一括して扱うための「母関数（特性多項式）」です。
* $T(x) = x(x-1)\prod (x-t_m)$ は、微分方程式の「確定特異点（動かない特異点 $0, 1$ と、動くパラメータ $t_m$）」を根に持つ多項式です。
* $L(x) = \prod (x-\lambda_k)$ は、「見かけの特異点」を根に持つ多項式です。
このように特異点の位置を多項式の根としてパックすることで、ガルニエ系全体の挙動を「$T(x)$ と $L(x)$ が織りなす1変数の代数関数論」として驚くほど見通し良く記述できるようになります。これは岡本和夫先生の研究の真骨頂とも言える手法です。
\end{Q}

\begin{Q}[【核心】$Z_l(t_h)$ に $x=t_h$ を代入するだけで、なぜ逆行列の証明が終わるのか？]
\textbf{A.} これは部分分数分解と「極（分母が0になる点）」の性質を巧みに利用したマジックです。
有理関数 $Z_l(x)$ の定義式を見ると、分子に $T(x) \propto (x-t_1)\cdots(x-t_g)$ があります。したがって、$h \neq l$ のときは単に $T(t_h) = 0$ となり $Z_l(t_h) = 0$ になります。
一方、$h = l$ のときは、分母にも $(x-t_l)$ があるため $0/0$ の不定形になります。ロピタルの定理のように約分して極限をとることで、分子は微分された $T'(t_l)$ となり、$0$ ではない値が残ります。
これを、部分分数展開した側（$\sum$ の形）と比較します。すると、「行列 $E$ と、$M_l M^{k,l}$ からなる行列 $F$ の積」が、「$h=l$ のときは $1$、$h \neq l$ のときは $0$」すなわち単位行列（クロネッカーのデルタ $\delta_{hl}$）になることが、一切の泥臭い連立方程式を解くことなく、極限の評価だけで鮮やかに証明されるのです。
\end{Q}

前節までの議論により、ハミルトニアン $H_h$ は以下の連立1次方程式を満たすことがわかっている。
\begin{equation*}
    \sum_{h=1}^g E_{kh} H_h = \mu_k^2 + \hat{a}_k \mu_k + \hat{b}_k
\end{equation*}
ここで $\hat{a}_k, \hat{b}_k$ は $H_h$ を含まない項であり、具体的には次で与えられる。
\begin{align*}
    \hat{a}_k &= \frac{1-\kappa_0}{\lambda_k} + \frac{1-\kappa_1}{\lambda_k-1} + \sum_{h=1}^g \frac{1-\theta_h}{\lambda_k-t_h} - \sum_{p=1 (p \neq k)}^g \frac{1}{\lambda_k-\lambda_p} \\
    \hat{b}_k &= \frac{\kappa}{\lambda_k(\lambda_k-1)} - \frac{2\lambda_k-1}{\lambda_k(\lambda_k-1)}\mu_k + \sum_{p=1 (p \neq k)}^g \frac{\lambda_p(\lambda_p-1)\mu_p}{\lambda_k(\lambda_k-1)(\lambda_k-\lambda_p)}
\end{align*}
行列 $E$ の逆行列 $F_{lk} = M_l M^{k,l}$ を左から掛けることで、$H_l$ は次のように表される。
\begin{equation}
    H_l = M_l \sum_{k=1}^g M^{k,l} \left( \mu_k^2 + \hat{a}_k \mu_k + \hat{b}_k \right) \label{eq:Hl_base}
\end{equation}

\begin{theorem}[ハミルトニアンの最終形態]
式 \eqref{eq:Hl_base} を $\mu_k$ について整理すると、以下の式 (3.73) が得られる。
\begin{equation}
    H_l = M_l \left[ \sum_{k=1}^g M^{k,l} \left\{ \mu_k^2 - A_{k,l}\mu_k + \frac{\kappa}{\lambda_k(\lambda_k-1)} \right\} \right]
\end{equation}
\begin{equation*}
    A_{k,l} = \frac{\kappa_0}{\lambda_k} + \frac{\kappa_1}{\lambda_k-1} + \sum_{h=1}^g \frac{\theta_h - \delta_{hl}}{\lambda_k - t_h}
\end{equation*}
\end{theorem}

\begin{proof}[証明と係数比較]
式 \eqref{eq:Hl_base} の括弧内を展開し、$\mu_k^2$ の項、$\mu_k^0$（定数）の項、$\mu_k^1$ の項に分けて比較する。
\begin{itemize}
    \item \textbf{2次の項}: そのまま $\mu_k^2$ となり一致する。
    \item \textbf{定数項}: $\hat{b}_k$ の第1項 $\frac{\kappa}{\lambda_k(\lambda_k-1)}$ がそのまま現れるため一致する。
    \item \textbf{1次の項}: ここが非自明である。$\hat{b}_k$ の二重和の部分は、添字 $k$ と $p$ を入れ替えることで次のように変形できる。
    \begin{equation*}
        \sum_{k=1}^g M^{k,l} \sum_{p \neq k} \frac{\lambda_p(\lambda_p-1)\mu_p}{\lambda_k(\lambda_k-1)(\lambda_k-\lambda_p)} = \sum_{k=1}^g \mu_k \sum_{p \neq k} M^{p,l} \frac{\lambda_k(\lambda_k-1)}{\lambda_p(\lambda_p-1)(\lambda_p-\lambda_k)}
    \end{equation*}
    したがって、式 \eqref{eq:Hl_base} における $\mu_k$ の全体の係数（$M^{k,l}$ をくくり出したもの）を $C_k$ とすると、
    \begin{equation}
        M^{k,l} C_k = M^{k,l} \left( \hat{a}_k - \frac{2\lambda_k-1}{\lambda_k(\lambda_k-1)} \right) + \sum_{p \neq k} M^{p,l} \frac{\lambda_k(\lambda_k-1)}{\lambda_p(\lambda_p-1)(\lambda_p-\lambda_k)}
    \end{equation}
    となる。この $C_k$ がちょうど $-A_{k,l}$ に等しいことを示せばよい。
\end{itemize}
代数的な計算により、$\hat{a}_k - \frac{2\lambda_k-1}{\lambda_k(\lambda_k-1)} = - A_{k,l} + \sum_{h \neq l} \frac{1}{\lambda_k-t_h} - \sum_{p \neq k} \frac{1}{\lambda_k-\lambda_p}$ がわかる。これを代入し、$M^{k,l}C_k = M^{k,l}(-A_{k,l})$ を示すことは、以下の等式を示すことと同値である。
\begin{equation}
    \sum_{h \neq l} \frac{1}{\lambda_k-t_h} - \sum_{p \neq k} \frac{1}{\lambda_k-\lambda_p} = \sum_{p \neq k} \frac{M^{p,l} \lambda_k(\lambda_k-1)}{M^{k,l} \lambda_p(\lambda_p-1)(\lambda_k-\lambda_p)} \label{eq:identity_target}
\end{equation}
多項式 $P(x) = \prod_{m \neq l}(x-t_m)$ および $L(x) = \prod (x-\lambda_q)$ を用いると、$\frac{M^{k,l}}{\lambda_k(\lambda_k-1)} = \frac{P(\lambda_k)}{L'(\lambda_k)}$ と書けるため、式 \eqref{eq:identity_target} の右辺は
\begin{equation*}
    \sum_{p \neq k} \frac{P(\lambda_p)L'(\lambda_k)}{P(\lambda_k)L'(\lambda_p)(\lambda_k-\lambda_p)}
\end{equation*}
となる。これは、有理関数の部分分数展開（ラグランジュ補間の微分の公式）から得られる恒等式と完全に一致する。よって $C_k = -A_{k,l}$ が証明され、式 (3.73) が導出された。
\end{proof}

\begin{Q}[【核心】なぜ二重和の計算で「添字 $k$ と $p$ の入れ替え」を行う必要があるのか？]
\textbf{A.} 行列の積を計算して $H_l = \sum_k F_{lk} (\cdots)$ と書き下した段階では、和の主役は $k$ です。しかし、$\hat{b}_k$ の中にさらに別の和 $\sum_{p \neq k} \mu_p$ が含まれているため、このままでは $\mu_1, \mu_2, \dots$ がそれぞれの項に散らばってしまいます。
最終目標である式 (3.73) は「各 $k$ について $\mu_k^2$ と $\mu_k^1$ を綺麗にまとめた形」をしています。したがって、散らばった $\mu$ を主役（係数）としてまとめ直すために、和の順番を $\sum_k \sum_p$ から $\sum_p \sum_k$ へと入れ替え、形式的に $p$ を $k$ に読み替えることで、全体の $\mu_k$ の係数を一つに集約させているのです。
\end{Q}

\begin{Q}[【深掘り】証明の最後に登場した「ラグランジュ補間の微分の公式」とは何か？]
\textbf{A.} これは代数幾何や可積分系において、複雑な有理関数の和を一瞬で計算するための強力な恒等式です。
一般に、$g-1$ 次の多項式 $P(x)$ は、ラグランジュ補間により $P(x) = \sum_{p=1}^g \frac{P(\lambda_p)}{L'(\lambda_p)} \frac{L(x)}{x-\lambda_p}$ と展開できます。この両辺を $x$ で微分し、$x = \lambda_k$ を代入して整理すると、
\begin{equation*}
    \frac{P'(\lambda_k)}{P(\lambda_k)} - \frac{L''(\lambda_k)}{2L'(\lambda_k)} = \sum_{p \neq k} \frac{P(\lambda_p)L'(\lambda_k)}{P(\lambda_k)L'(\lambda_p)(\lambda_k-\lambda_p)}
\end{equation*}
という美しい関係式が自動的に得られます。左辺の対数微分 $\frac{P'}{P}$ や $\frac{L''}{L'}$ を計算すると、まさに式 \eqref{eq:identity_target} の左辺のシグマの差そのものになります。パンルヴェ理論のハミルトニアンがこれほど綺麗にまとまる背後には、こうした代数関数論の深い恒等式が隠れているのです。
\end{Q}

\begin{story}

以下，前節から考察している2階フックス型線型常微分方程式
\begin{equation}
    \frac{d^2 y}{dx^2} + p_1(x,t) \frac{dy}{dx} + p_2(x,t)y = 0
\end{equation}
のモノドロミー保存変形の計算に戻ろう．係数は(3.67), (3.68)により与えられていた．我々は，$x=\lambda_k$ が見かけの特異点である，という仮定を使って，$H_l$ を他の変数の有理関数として具体的に求めたところである．

ここで，第3.2節において調べたことを思いだそう．(3.33)に対して
\begin{equation}
    r(x,t) = -p_2(x,t) + \frac{1}{4}p_1(x,t)^2 + \frac{1}{2}\frac{\partial}{\partial x}p_1(x,t)
\end{equation}
とし，SL型方程式
\begin{equation}
    \frac{d^2 z}{dx^2} = r(x,t)z
\end{equation}
を併せて考える．ここで124ページの考察を補題の形にまとめておこう．
\end{story}

124ページの命題3.7では
\[
\frac{d\phi}{dx}+\frac{1}{2}p_{1}(x,t)\phi=0
\]
のモノドロミーが$t$によらないことを仮定していたが、すぐ後で$\mathbb{P}^{1}(\mathbb{C})$では自動的にその条件が満たされることを言っている。具体的に解くと、
\[
\phi(x,t)=\prod_{i=1}^{M}(x-\xi_{i})^{-\frac{1}{2}\alpha_{i}}
\]
で与えられてモノドロミーは$e^{2\pi i\alpha_{i}(t)}$というスカラーで書けている。
\[
\frac{\partial^{2}\vec{\phi}}{\partial x^{2}} + p_{1}(x,t)\frac{\partial\vec{\phi}}{\partial x}+p_{2}(x,t)\vec{\phi}
\]
のモノドロミーが$t$によらないならモノドロミー行列は対角行列で、その成分は、$x=\xi$での特性指数を$s_{1},s_{2}$とすると、$e^{2\pi is_{i}}$の形であるから$s_{i}$が$t$によらないことは必要。つまり、$\alpha_{i}$は特性指数の和で書けているので$t$によらず、$\phi$の$x=\xi_{i}$におけるモノドロミー行列は$t$によらない。よって以下が導ける。

\begin{story}

\begin{lemma}[3.5]
(3.33)がモノドロミー保存変形を定めることと，(3.25)がそうなることは同値であり，さらにこれは，以下の偏微分方程式系が $x$ の有理関数を解としてもつことと同値である．

\begin{align}
    \frac{\partial^3}{\partial x^3}w_l - 4r(x,t)\frac{\partial}{\partial x}w_l - 2\frac{\partial r}{\partial x}(x,t)w_l + 2\frac{\partial r}{\partial t_l}(x,t) = 0 \label{eq:3.82} \\
    \left( \frac{\partial}{\partial t_h} - w_h \frac{\partial}{\partial x} \right) w_l = \left( \frac{\partial}{\partial t_l} - w_l \frac{\partial}{\partial x} \right) w_h \quad (h,l = 1,\cdots,g) \label{eq:3.83}
\end{align}

\end{lemma}
\end{story}

\begin{Q}[【深掘り】なぜわざわざ元の微分方程式 (3.33) を、1階微分の項がないSL型 (3.25) に変換して考える必要があるのか？]
\textbf{A.} モノドロミー保存変形の条件（Lax方程式などの非線形偏微分方程式系）を導出する際、**1階微分の項 $\frac{dy}{dx}$ が存在すると計算が絶望的に複雑になるから**です。
実は、未知関数 $y$ に対して $z = y \exp(\frac{1}{2}\int p_1 dx)$ という変換（ゲージ変換）を行うと、1階微分の項が完全に消去され、$z'' = r(x,t)z$ というシンプルな形（Sturm-Liouville型、あるいはシュレディンガー方程式型）に帰着させることができます。
この変換は、解の「特異点周りでの回り方（モノドロミー表現）」の本質を一切変えません。そのため、「モノドロミーが保存される」という幾何学的な条件を計算するときは、見通しの良いSL型 (3.25) を使って変形方程式を書き下すのが、可積分系理論における絶対的な定石なのです。(3.32) の $r(x,t)$ の式は、まさにそのゲージ変換から生じるポテンシャル項を表しています。
\end{Q}

\begin{story}
実際，命題3.8により(3.82), (3.83)の有理関数解 $w_l = a_2^{(l)}(x,t)$ は存在すれば唯一であり，それに対して $a_1^{(l)}(x,t)$ を(3.23)により定めると，変形方程式系
\begin{equation}
\left\{
\begin{array}{l}
    \displaystyle \frac{\partial^2 \vec{\phi}}{\partial x^2} + p_1(x,t)\frac{\partial \vec{\phi}}{\partial x} + p_2(x,t)\vec{\phi} = 0 \vspace{0.5em} \\
    \displaystyle \frac{\partial \vec{\phi}}{\partial t_l} = a_1^{(l)}(x,t)\vec{\phi} + a_2^{(l)}(x,t)\frac{\partial \vec{\phi}}{\partial x}
\end{array}
\right. \label{eq:3.20}
\end{equation}
は完全積分可能である．さらに次の命題も成り立つ．

\begin{proposition}[3.14]
偏微分方程式系(3.82)の有理関数解 $w_l$ は，(3.83)を満たす．
\end{proposition}

この命題により，モノドロミー保存変形は(3.82)の考察に帰着される．証明のために
\begin{equation*}
    L = \frac{\partial^3}{\partial x^3} - 4r(x,t)\frac{\partial}{\partial x} - 2\frac{\partial r}{\partial x}(x,t)
\end{equation*}
とおく．計算によって次の補題の成立を確認することができる．

\begin{lemma}[3.6]
任意の解析関数 $f,g$ に対して次式が成り立つ．
\begin{equation*}
    2\frac{\partial f}{\partial x}L(g) - 2\frac{\partial g}{\partial x}L(f) + f\frac{\partial}{\partial x}L(g) - g\frac{\partial}{\partial x}L(f) = L\left( f\frac{\partial g}{\partial x} - g\frac{\partial f}{\partial x} \right)
\end{equation*}
\end{lemma}

\end{story}

\begin{proof}
\begin{align*}
    & 2 \frac{\partial f}{\partial x} L(g) - 2 \frac{\partial g}{\partial x} L(f) + f \frac{\partial}{\partial x} L(g) - g \frac{\partial}{\partial x} L(f) \\
    &= 2 \frac{\partial f}{\partial x} \left( \frac{\partial^3 g}{\partial x^3} - 4r \frac{\partial g}{\partial x} - 2 \frac{\partial r}{\partial x} g \right) - 2 \frac{\partial g}{\partial x} \left( \frac{\partial^3 f}{\partial x^3} - 4r \frac{\partial f}{\partial x} - 2 \frac{\partial r}{\partial x} f \right) \\
    &\quad + f \frac{\partial}{\partial x} \left( \frac{\partial^3 g}{\partial x^3} - 4r \frac{\partial g}{\partial x} - 2 \frac{\partial r}{\partial x} g \right) - g \frac{\partial}{\partial x} \left( \frac{\partial^3 f}{\partial x^3} - 4r \frac{\partial f}{\partial x} - 2 \frac{\partial r}{\partial x} f \right) \\
    &= 2 \left( \frac{\partial f}{\partial x}\frac{\partial^3 g}{\partial x^3} - \frac{\partial g}{\partial x}\frac{\partial^3 f}{\partial x^3} \right) - 4 \frac{\partial r}{\partial x} \left( \frac{\partial f}{\partial x}g - \frac{\partial g}{\partial x}f \right) \\
    &\quad + f \left( \frac{\partial^4 g}{\partial x^4} - 4 \frac{\partial r}{\partial x}\frac{\partial g}{\partial x} - 4r \frac{\partial^2 g}{\partial x^2} - 2 \frac{\partial^2 r}{\partial x^2}g - 2 \frac{\partial r}{\partial x}\frac{\partial g}{\partial x} \right) \\
    &\quad - g \left( \frac{\partial^4 f}{\partial x^4} - 4 \frac{\partial r}{\partial x}\frac{\partial f}{\partial x} - 4r \frac{\partial^2 f}{\partial x^2} - 2 \frac{\partial^2 r}{\partial x^2}f - 2 \frac{\partial r}{\partial x}\frac{\partial f}{\partial x} \right) \\
    &= \left( 2\frac{\partial f}{\partial x}\frac{\partial^3 g}{\partial x^3} + f\frac{\partial^4 g}{\partial x^4} - 2\frac{\partial g}{\partial x}\frac{\partial^3 f}{\partial x^3} - g\frac{\partial^4 f}{\partial x^4} \right) - 4r \left( f\frac{\partial^2 g}{\partial x^2} - g\frac{\partial^2 f}{\partial x^2} \right) - 2\frac{\partial r}{\partial x} \left( f\frac{\partial g}{\partial x} - g\frac{\partial f}{\partial x} \right)
\end{align*}

ここで、
\begin{align*}
    \frac{\partial^3}{\partial x^3} \left( f\frac{\partial g}{\partial x} - g\frac{\partial f}{\partial x} \right) &= \frac{\partial^2}{\partial x^2} \left( f\frac{\partial^2 g}{\partial x^2} - g\frac{\partial^2 f}{\partial x^2} \right) \\
    &= \frac{\partial}{\partial x} \left( \frac{\partial f}{\partial x}\frac{\partial^2 g}{\partial x^2} + f\frac{\partial^3 g}{\partial x^3} - \frac{\partial g}{\partial x}\frac{\partial^2 f}{\partial x^2} - g\frac{\partial^3 f}{\partial x^3} \right) \\
    &= \left( 2\frac{\partial f}{\partial x}\frac{\partial^3 g}{\partial x^3} + f\frac{\partial^4 g}{\partial x^4} - 2\frac{\partial g}{\partial x}\frac{\partial^3 f}{\partial x^3} - g\frac{\partial^4 f}{\partial x^4} \right)
\end{align*}

よって、
\begin{align*}
    (\text{与式}) &= \frac{\partial^3}{\partial x^3} \left( f\frac{\partial g}{\partial x} - g\frac{\partial f}{\partial x} \right) - 4r \frac{\partial}{\partial x} \left( f\frac{\partial g}{\partial x} - g\frac{\partial f}{\partial x} \right) - 2\frac{\partial r}{\partial x} \left( f\frac{\partial g}{\partial x} - g\frac{\partial f}{\partial x} \right) \\
    &= L \left( f\frac{\partial g}{\partial x} - g\frac{\partial f}{\partial x} \right)
\end{align*}
\end{proof}

\begin{story}

\begin{proof}[命題3.14の証明]
(3.82)に有理関数解 $w_l$ が存在するとして
\begin{equation}\label{eq_3.6:L_r_wl}
    L(w_l) + 2\frac{\partial r}{\partial t_l}(x,t) = 0
\end{equation}
から，積分可能条件として
\[
\frac{\partial^{2} r}{\partial t_l\partial t_{h}}(x,t)=\frac{\partial^{2} r}{\partial t_h\partial t_{l}}(x,t)
\]
より、($w$の交換関係は当たり前に仮定していて、そこからどういう数式が生まれるか見ている。)
\begin{align}
    0 &= \frac{\partial}{\partial t_h}L(w_l) - \frac{\partial}{\partial t_l}L(w_h)\notag\\
    &= \left(\frac{\partial^{4}}{\partial t_h\partial x^{3}}-4\frac{\partial r}{\partial t_{h}}\frac{\partial}{\partial x}-4r\frac{\partial^{2}}{\partial t_{h}\partial x}-2\frac{\partial^{2}r}{\partial t_{h}\partial x}-2\frac{\partial r}{\partial t_{h}}\frac{\partial}{\partial x}\right)w_{l}\notag\\
    &-\left(\frac{\partial^{4}}{\partial t_l\partial x^{3}}-4\frac{\partial r}{\partial t_{l}}\frac{\partial}{\partial x}-4r\frac{\partial^{2}}{\partial t_{l}\partial x}-2\frac{\partial^{2}r}{\partial t_{l}\partial x}-2\frac{\partial r}{\partial t_{l}}\frac{\partial}{\partial x}\right)w_{h}\notag\\
    &= \left(\frac{\partial^{3}}{\partial x^{3}}-4r\frac{\partial}{\partial x}-2\frac{\partial r}{\partial x}\right)(\frac{\partial w_{l}}{\partial t_{h}}-\frac{\partial w_{h}}{\partial t_{l}})\notag\\
    &+\left(-4\frac{\partial r}{\partial t_h}\frac{\partial w_l}{\partial x} + 4\frac{\partial r}{\partial t_l}\frac{\partial w_h}{\partial x} - 2w_l\frac{\partial^2 r}{\partial x \partial t_h} + 2w_h\frac{\partial^2 r}{\partial x \partial t_l}\right)\notag\\
    &= L\left( \frac{\partial w_l}{\partial t_h} - \frac{\partial w_h}{\partial t_l} \right) + L_{hl} \label{eq:3.84}
\end{align}
が成り立っている．ここで
\begin{equation*}
    L_{hl} = -4\frac{\partial r}{\partial t_h}\frac{\partial w_l}{\partial x} + 4\frac{\partial r}{\partial t_l}\frac{\partial w_h}{\partial x} - 2w_l\frac{\partial^2 r}{\partial x \partial t_h} + 2w_h\frac{\partial^2 r}{\partial x \partial t_l}
\end{equation*}
であるが、\eqref{eq_3.6:L_r_wl}より、
\begin{align*}
    L_{hl} &=2\frac{\partial w_l}{\partial x}L(w_h) - 2\frac{\partial w_h}{\partial x}L(w_l)+w_{l}\frac{\partial}{\partial x}\left(-2\frac{\partial r}{\partial t_{h}}\right)-w_{h}\frac{\partial}{\partial x}\left(-2\frac{\partial r}{\partial t_{l}}\right)\\
    &= 2\frac{\partial w_l}{\partial x}L(w_h) - 2\frac{\partial w_h}{\partial x}L(w_l) + w_l\frac{\partial}{\partial x}L(w_h) - w_h\frac{\partial}{\partial x}L(w_l) \\
\end{align*}
補題3.6を使うと
\begin{equation*}
    L_{hl} = L\left( w_l\frac{\partial w_h}{\partial x} - w_h\frac{\partial w_l}{\partial x} \right)
\end{equation*}
したがって(3.84)より
\begin{equation*}
    L\left( \frac{\partial w_l}{\partial t_h} - \frac{\partial w_h}{\partial t_l} + w_l\frac{\partial w_h}{\partial x} - w_h\frac{\partial w_l}{\partial x} \right) = 0
\end{equation*}
となる．命題3.8の証明で見たように，仮定(P2), (P3)によれば，3階線型常微分方程式 $L(w)=0$ は非自明な有理関数解をもたない．(以下を参照)すなわち
\begin{equation*}
    \frac{\partial w_l}{\partial t_h} - \frac{\partial w_h}{\partial t_l} + w_l\frac{\partial w_h}{\partial x} - w_h\frac{\partial w_l}{\partial x} = 0
\end{equation*}
以上の議論は，(3.83)が(3.82)の積分可能条件に他ならないことを示している．
\end{proof}
\end{story}

\begin{Q}[【核心】作用素 $L$ と方程式(3.82)は、どこかで見たことがある形をしているが？]
\textbf{A.} お気付きの通り、以前のノートにあった**KdV方程式のLax対の生成子 $M_3$ と本質的に全く同じもの**です。
ノートでは $M_3 = -4\partial_x^3 - 6u\partial_x - 3u_x$ でした。これを定数倍でスケール調整し、$u$ をポテンシャル $r(x,t)$ に読み替えると、まさに $L = \partial_x^3 - 4r\partial_x - 2r_x$ になります。
シュレディンガー作用素 $\partial_x^2 - r(x,t)$ の固有値（あるいはモノドロミー）を保存するような変形を考えると、必ずこの3階の微分作用素 $L$ が現れます。(3.82) は $L(w_l) = -2\partial_{t_l}r$ と書けますが、これはソリトン理論における「Lenard関係式」や「KdV階層」を定義する方程式そのものです。モノドロミー保存変形（ガルニエ系）とは、本質的に「有限自由度のKdV階層」に他ならないことが、この式から明確に読み取れます。
\end{Q}

\begin{Q}[【深掘り】命題3.14は、なぜそれほどまでに重要なのか？]
\textbf{A.} モノドロミー保存変形を成立させるためには、本来 **(3.82)（ポテンシャルの時間発展）と (3.83)（異なる時間 $t_l, t_h$ のフローの可換性）の両方を同時に満たす $w_l$** を探さなければなりません。複数の時間を扱う多変数系（ガルニエ系）において、(3.83)のような非線形の連立微分方程式を解くことは通常至難の業です。
しかし命題3.14は、**「(3.82)の有理関数解を見つけさえすれば、面倒な(3.83)は自動的に満たされる」**という驚くべき事実を保証しています。つまり、問題を解くための労力が半分（あるいはそれ以下）に劇的に減るのです。
この「奇跡」を代数的に証明するための鍵となるのが、補題3.6の恒等式です。作用素 $L$ が持つ美しい代数的な対称性のおかげで、個別の時間発展(3.82)が保証されれば、系全体の無矛盾性(3.83)も自動的に担保される構造になっています。
\end{Q}

\begin{Q}[【核心】作用素 $L$ で括った中身が $=0$ になるのはなぜか？積分定数のようなものは出ないのか？]
\textbf{A.} ここがこの証明の最大の要です。一般の微分方程式 $L(W) = 0$ には当然、非自明な解（積分定数を含む関数の族）が存在します。しかしここで探しているのは「有理関数解」です。
パンルヴェ理論の枠組みでは、特異点の配置に関する一般的な仮定（P2, P3など）を置くことで、「$L(w)=0$ を満たす有理関数は $w=0$ しか存在しない（核が自明である）」という事実が代数幾何的に保証されています。
非線形偏微分方程式の可換性条件 (3.83) という極めて複雑な等式を、「線形作用素 $L$ の核が自明である」という代数的な性質に帰着させて証明し切ってしまう、極めて美しい論法です。
\end{Q}

\begin{story}
そこで，(3.32)を用いて $r(x,t)$ を計算しよう．実際，$r(x,t)$ は $x=0, x=1, x=t$ で高々2位の極をもつ有理関数である．特性指数を考慮すれば，$r(x,t)$ は

\end{story}

\begin{story}
これは、$x=0,1,t_{l}$で確定特異点を持つので$p_{2}$が高々2位で、$p_{1}$が高々1位なので、$p_{1}^{2},\frac{\partial p_{1}}{\partial x}$が高々2位の極を持つことから従う。
$r(x,t)$ は

\begin{align}
    r(x,t) &= \frac{\alpha_0^{2}}{x^2} + \frac{\alpha_1^{2}}{(x-1)^2} +\sum_{l=1}^g \frac{\beta_l^{2}}{(x-t_l)^2}+\sum_{k=1}^g\frac{\gamma_{k}^{2}}{(x-\lambda_k)^2} \nonumber\\
    &+\frac{\alpha_{0}^{1}}{x} + \frac{\alpha_{1}^{1}}{x-1} + \sum_{l=1}^g  \frac{\beta_l^{1}}{x-t_l}+ \sum_{k=1}^g \frac{\gamma_{k}^{1}}{x-\lambda_k} 
\end{align}

という形をしている。ここで、$x=\infty$で確定特異点を持つことから、命題2.15より、
\[
\frac{1}{u^{2}}r\left(\frac{1}{u},t\right)
\]
が$u=0$で正則である必要がある。ここで主要部は
\begin{align*}
    & \frac{1}{u^2} \left( u\alpha_0^{1} + \frac{\alpha_1^{1}}{1/u - 1} + \sum_{l=1}^g \frac{\beta_l^{1}}{1/u - t_l} + \sum_{k=1}^g \frac{\gamma_k^{1}}{1/u - \lambda_k} \right) \\
    &= \frac{\alpha_0^{1}}{u} + \frac{\alpha_1^{1}}{u(1-u)} + \sum_{l=1}^g \frac{\beta_l^{1}}{u(1-t_l u)} + \sum_{k=1}^g \frac{\gamma_k^{1}}{u(1-\lambda_k u)}
\end{align*}
よって、$u=0$ で 正則となるには、
\begin{equation*}
    \alpha_0^{1} + \alpha_1^{1} + \sum_{l=1}^g \beta_l^{1} + \sum_{k=1}^g \gamma_k^{1} = 0
\end{equation*}
が必要。このとき、$\alpha_0^{1}$ を消去して整理すると、
\begin{align*}
    & \frac{1}{x} \left( -\alpha_1^{1} - \sum_{l=1}^g \beta_l^{1} - \sum_{k=1}^g \gamma_k^{1} \right) + \frac{\alpha_1^{1}}{x-1} + \sum_{l=1}^g \frac{\beta_l^{1}}{x-t_l} + \sum_{k=1}^g \frac{\gamma_k^{1}}{x-\lambda_k} \\
    &= \frac{\alpha_1^{1}}{x(x-1)} + \sum_{l=1}^g \frac{t_l \beta_l^{1}}{x(x-t_l)} + \sum_{k=1}^g \frac{\lambda_k \gamma_k^{1}}{x(x-\lambda_k)}
\end{align*}
$\alpha_1^{1}, \beta_l^1, \gamma_k^1$ を適当につけ替えることで、
\begin{equation*}
    \frac{\alpha_1^1}{x(x-1)} + \sum_{l=1}^g \frac{t_l(t_l-1)\beta_l^1}{x(x-1)(x-t_l)} - \sum_{k=1}^g \frac{\lambda_k(\lambda_k-1)\gamma_k^1}{x(x-1)(x-\lambda_k)}
\end{equation*}
となる。よって、$\alpha_1^1 \to \alpha_\infty, \ \beta_l^1 \to K_l, \ \gamma_k^1 \to \nu_k, \ \alpha_0^2 \to \alpha_0, \ \alpha_1^2 \to \alpha_1, \ \beta_l^2 \to \beta_l, \ \gamma_k^2 \to \gamma_k$ として、

\begin{align}
    r(x,t) &= \frac{\alpha_0}{x^2} + \frac{\alpha_1}{(x-1)^2} + \frac{\alpha_\infty}{x(x-1)} \nonumber \\
    &\quad + \sum_{l=1}^g \left[ \frac{\beta_l}{(x-t_l)^2} + \frac{t_l(t_l-1)K_l}{x(x-1)(x-t_l)} \right] \label{eq:3.85} \\
    &\quad + \sum_{k=1}^g \left[ \frac{3}{4(x-\lambda_k)^2} - \frac{\lambda_k(\lambda_k-1)\nu_k}{x(x-1)(x-\lambda_k)} \right] \nonumber
\end{align}

という形をしている．他方，(3.32)に(3.67)と(3.68)を代入して，係数の関係を調べる．すると

変換後のリーマン図式が
\begin{equation*}
\left\{
\begin{array}{ccccc}
x=0 & x=1 & x=t_l & x=\lambda_k & x=\infty \\
\frac{1+\kappa_0}{2} & \frac{1+\kappa_1}{2} & \frac{1+\theta_l}{2} & \frac{3}{2} & \frac{1+\kappa_\infty}{2} \\
\frac{1-\kappa_0}{2} & \frac{1-\kappa_1}{2} & \frac{1-\theta_l}{2} & -\frac{1}{2} & \frac{1-\kappa_\infty}{2}
\end{array}
\right\}
\end{equation*}
となることと、
\begin{equation}
    \frac{d^2 z}{dx^2} - r(x,t)z = 0
\end{equation}
から、

\begin{equation*}
    \alpha_0 = \frac{1}{4}(\kappa_0^2 - 1), \quad \alpha_1 = \frac{1}{4}(\kappa_1^2 - 1), \quad \beta_l = \frac{1}{4}(\theta_l^2 - 1)
\end{equation*}

また，$\frac{d^2 z}{dx^2} - r(x,t)z = 0$ の決定方程式は命題2.15から，
\begin{equation*}
    -s(-s-1) - \left. \frac{1}{u^2} r\left(\frac{1}{u}, t\right) \right|_{u=0} = 0
\end{equation*}
より，
\begin{align*}
    \frac{1-\kappa_\infty^2}{4} &= - \left( \frac{1}{u^2} \left( u^2\alpha_0 + \frac{\alpha_1 u^2}{(1-u)^2} + \frac{\alpha_\infty u^2}{1-u} + \sum_{l=1}^g \left[ \frac{\beta_l u^2}{(1-t_l u)^2} + \frac{u^3 t_l(t_l-1)K_l}{(1-u)(1-t_l u)} \right] \right. \right. \\
    &\qquad \left. \left. + \sum_{k=1}^g \left[ \frac{\gamma_k u^2}{(1-\lambda_k u)^2} - \frac{u^3 \lambda_k(\lambda_k-1)\nu_k}{(1-u)(1-\lambda_k u)} \right] \right) \right) \Bigg|_{u=0} \\
    &= - \left( \alpha_0 + \frac{\alpha_1}{(1-u)^2} + \frac{\alpha_\infty}{1-u} + \sum_{l=1}^g \frac{\beta_l}{(1-t_l u)^2} + \sum_{k=1}^g \frac{\gamma_k}{(1-\lambda_k u)^2} \right) \Bigg|_{u=0} \\
    &= -\left( \alpha_0 + \alpha_1 + \alpha_\infty + \sum_{l=1}^g \beta_l + \sum_{k=1}^g \gamma_k \right)
\end{align*}
整理して $\alpha_\infty$ について解くと、
\begin{align*}
    \alpha_\infty &= \frac{1-\kappa_0^2}{4} + \frac{1-\kappa_1^2}{4} + \sum_{l=1}^g \frac{1-\theta_l^2}{4} - \sum_{k=1}^g \frac{3}{4} - \frac{1-\kappa_\infty^2}{4} \\
    &= -\frac{1}{4} \left( \kappa_0^2 + \kappa_1^2 + \sum_{l=1}^g \theta_l^2 - \kappa_\infty^2 - 1 \right) - \frac{g}{2}
\end{align*}
\end{story}

\begin{Q}[【深掘り】式(3.85)の見かけの特異点 $x=\lambda_k$ の項にある「$3/4$」という数字はどこから来たのか？]
\textbf{A.} これはSL型方程式（シュレディンガー方程式型）における**見かけの特異点の絶対的な刻印**です。
フックス型微分方程式をSL型にゲージ変換（$z = y \exp(\frac{1}{2}\int p_1 dx)$）すると、変換後のポテンシャル $r(x,t)$ の各特異点での2位の極の留数は、特性指数の差 $\Delta \rho$ を用いて $\frac{1}{4}((\Delta \rho)^2 - 1)$ という普遍的な形で書けることが知られています（シュワルツ微分の性質）。
確定特異点 $x=0, 1, t_l$ では特性指数の差がそれぞれ $\kappa_0, \kappa_1, \theta_l$ なので、$\alpha_0, \alpha_1, \beta_l$ はまさにこの公式通りになっています。
そして、見かけの特異点 $x=\lambda_k$ では、以前確認したように「特性指数は $0, 2$」です。したがってその差は $\Delta \rho = 2 - 0 = 2$ となります。これを公式に代入すると、
$$ \frac{1}{4}(2^2 - 1) = \frac{3}{4} $$
となります。つまりこの $\frac{3}{4}$ は、そこが「特性指数の差が $2$ である見かけの特異点である」という幾何学的な事実から代数的に必然として導かれる定数なのです。
\end{Q}

\begin{story}
また、$r(x,t) = -p_2 + \frac{1}{4}p_1^2 + \frac{1}{2}\frac{\partial p_1}{\partial x}$
を用いて、$x=t_l$ の1位の極に寄与するものを集めると、
\begin{equation*}
    \frac{t_l(t_l-1)H_l}{x(x-1)(x-t_l)} + \frac{1}{2} \frac{1-\theta_l}{x-t_l} \left( \frac{1-\kappa_0}{x} + \frac{1-\kappa_1}{x-1} + \sum_{m \neq l} \frac{1-\theta_m}{x-t_m} - \sum_{k=1}^g \frac{1}{x-\lambda_k} \right)
\end{equation*}
これと $\frac{t_l(t_l-1)K_l}{x(x-1)(x-t_l)}$ の $x=t_l$ での留数を比べることで、
\begin{equation*}
    K_l = H_l + \frac{1-\theta_l}{2} \left( \frac{1-\kappa_0}{t_l} + \frac{1-\kappa_1}{t_l-1} + \sum_{m \neq l} \frac{1-\theta_m}{t_l-t_m} - \sum_{k=1}^g \frac{1}{t_l-\lambda_k} \right)
\end{equation*}
よって、$W_l := \frac{1-\kappa_0}{t_l} + \frac{1-\kappa_1}{t_l-1} + \sum_{m \neq l} \frac{1-\theta_m}{t_l-t_m} + \sum_{k=1}^g \frac{1}{\lambda_k-t_l}$ とおけば、
\begin{equation*}
    K_l = H_l + \frac{1-\theta_l}{2} W_l
\end{equation*}

同様に、$x=\lambda_k$ の1位の極に寄与するものを集めると
\begin{equation*}
    -\frac{\lambda_k(\lambda_k-1)\mu_k}{x(x-1)(x-\lambda_k)} - \frac{1}{2} \frac{1}{x-\lambda_k} \left( \frac{1-\kappa_0}{x} + \frac{1-\kappa_1}{x-1} + \sum_{l=1}^g \frac{1-\theta_l}{x-t_l} - \sum_{m \neq k} \frac{1}{x-\lambda_m} \right)
\end{equation*}
これと $-\frac{\lambda_k(\lambda_k-1)\nu_k}{x(x-1)(x-\lambda_k)}$ の $x=\lambda_k$ での留数を比べることで、
\begin{equation*}
    \nu_k = \mu_k + \frac{1}{2} \left( \frac{1-\kappa_0}{\lambda_k} + \frac{1-\kappa_1}{\lambda_k-1} + \sum_{l=1}^g \frac{1-\theta_l}{\lambda_k-t_l} - \sum_{m \neq k} \frac{1}{\lambda_k-\lambda_m} \right)
\end{equation*}
よって、$\omega_k := \frac{1-\kappa_0}{\lambda_k} + \frac{1-\kappa_1}{\lambda_k-1} + \sum_{l=1}^g \frac{1-\theta_l}{\lambda_k-t_l} - \sum_{m \neq k} \frac{1}{\lambda_k-\lambda_m}$ とおけば、
\begin{equation*}
    \nu_k = \mu_k + \frac{1}{2} \omega_k
\end{equation*}
まとめると、
\begin{equation}
    K_l = H_l + \frac{1-\theta_l}{2} W_l, \quad \nu_k = \mu_k + \frac{1}{2} \omega_k \label{eq:3.86}
\end{equation}
となる．ただし，(3.86)では次のようにおいた．
\begin{align*}
    W_l &= \frac{1-\kappa_0}{t_l} + \frac{1-\kappa_1}{t_l-1} + \sum_{h=1 (h\neq l)}^g \frac{1-\theta_h}{t_l-t_h} + \sum_{k=1}^g \frac{1}{\lambda_k-t_l} \\
    \omega_k &= \frac{1-\kappa_0}{\lambda_k} + \frac{1-\kappa_1}{\lambda_k-1} + \sum_{l=1}^g \frac{1-\theta_l}{\lambda_k-t_l} - \sum_{p=1 (p\neq k)}^g \frac{1}{\lambda_k-\lambda_p}
\end{align*}

ここで(3.86)を，変換
\begin{equation}
    (\lambda, \mu, H, t) \mapsto (\lambda, \nu, K, t) \label{eq:3.87}
\end{equation}
を定義する，と見る．このとき
\begin{equation*}
    \left(\frac{\partial}{\partial \lambda_p}\right)\omega_k = \frac{\partial}{\partial \lambda_p}\left( -\frac{1}{\lambda_k-\lambda_p} \right) = -\frac{1}{(\lambda_p-\lambda_k)^2}
\end{equation*}
よって、
\begin{equation*}
    \left(\frac{\partial}{\partial \lambda_p}\right)\omega_k - \left(\frac{\partial}{\partial \lambda_k}\right)\omega_p = -\frac{1}{(\lambda_p-\lambda_k)^2} + \frac{1}{(\lambda_k-\lambda_p)^2} = 0
\end{equation*}


\begin{equation*}
    (1-\theta_h)\left(\frac{\partial}{\partial t_l}\right)W_h = (1-\theta_h)\frac{\partial}{\partial t_l}\left( \frac{1-\theta_l}{t_h-t_l} \right) = \frac{(1-\theta_h)(1-\theta_l)}{(t_h-t_l)^2}
\end{equation*}
よって、
\begin{align*}
    (1-\theta_h)\left(\frac{\partial}{\partial t_l}\right)W_h - (1-\theta_l)\left(\frac{\partial}{\partial t_h}\right)W_l &= \frac{(1-\theta_h)(1-\theta_l)}{(t_h-t_l)^2} - \frac{(1-\theta_l)(1-\theta_h)}{(t_l-t_h)^2} = 0
\end{align*}


\begin{align*}
    (1-\theta_l)\left(\frac{\partial}{\partial \lambda_k}\right)W_l &= (1-\theta_l)\frac{\partial}{\partial \lambda_k}\left( \frac{1}{\lambda_k-t_l} \right) = -\frac{1-\theta_l}{(\lambda_k-t_l)^2} \\
    \left(\frac{\partial}{\partial t_l}\right)\omega_k &= \frac{\partial}{\partial t_l}\left( \frac{1-\theta_l}{\lambda_k-t_l} \right) = \frac{1-\theta_l}{(\lambda_k-t_l)^2}
\end{align*}
よって、
\begin{equation*}
    (1-\theta_l)\left(\frac{\partial}{\partial \lambda_k}\right)W_l + \left(\frac{\partial}{\partial t_l}\right)\omega_k = -\frac{1-\theta_l}{(\lambda_k-t_l)^2} + \frac{1-\theta_l}{(\lambda_k-t_l)^2} = 0
\end{equation*}
より、
\begin{align*}
    \left(\frac{\partial}{\partial \lambda_p}\right)\omega_k - \left(\frac{\partial}{\partial \lambda_k}\right)\omega_p &= 0 \\
    (1-\theta_h)\left(\frac{\partial}{\partial t_l}\right)W_h - (1-\theta_l)\left(\frac{\partial}{\partial t_h}\right)W_l &= 0 \\
    (1-\theta_l)\left(\frac{\partial}{\partial \lambda_k}\right)W_l + \left(\frac{\partial}{\partial t_l}\right)\omega_k &= 0
\end{align*}
となることに注意すると
\begin{equation*}
    \sum_{k=1}^g d\mu_k \wedge d\lambda_k - \sum_{l=1}^g dH_l \wedge dt_l = \sum_{k=1}^g d\nu_k \wedge d\lambda_k - \sum_{l=1}^g dK_l \wedge dt_l
\end{equation*}
が成り立つことがわかる．ただし $\left(\frac{\partial}{\partial t_h}\right), \left(\frac{\partial}{\partial \lambda_p}\right)$ 等の記号は，$W_l$ や $\omega_k$ に陽に現れる変数 $t_h, \lambda_p$ についての偏微分を表す．

実際、\begin{align*}
    \sum_{k=1}^g d\nu_k \wedge d\lambda_k - \sum_{l=1}^g dK_l \wedge dt_l &= \sum_{k=1}^g d\left(\mu_k + \frac{1}{2}\omega_k\right) \wedge d\lambda_k - \sum_{l=1}^g d\left(H_l + \frac{1}{2}(1-\theta_l)W_l\right) \wedge dt_l \\
    &= \sum_{k=1}^g d\mu_k \wedge d\lambda_k - \sum_{l=1}^g dH_l \wedge dt_l + \left( \frac{1}{2}\sum_{k=1}^g d\omega_k \wedge d\lambda_k - \frac{1}{2}\sum_{l=1}^g (1-\theta_l)dW_l \wedge dt_l \right)
\end{align*}

ここで
\begin{align*}
    d\omega_k \wedge d\lambda_k &= \sum_{l=1}^g \frac{\partial \omega_k}{\partial t_l} dt_l \wedge d\lambda_k + \sum_{\substack{m=1 \\ (m\neq k)}}^g \frac{\partial \omega_k}{\partial \lambda_m} d\lambda_m \wedge d\lambda_k \\
    dW_l \wedge dt_l &= \sum_{\substack{m=1 \\ (m\neq l)}}^g \frac{\partial W_l}{\partial t_m} dt_m \wedge dt_l + \sum_{k=1}^g \frac{\partial W_l}{\partial \lambda_k} d\lambda_k \wedge dt_l
\end{align*}

より、
\begin{align*}
    & \frac{1}{2}\sum_{k=1}^g d\omega_k \wedge d\lambda_k - \frac{1}{2}\sum_{l=1}^g (1-\theta_l)dW_l \wedge dt_l \\
    &= \frac{1}{2}\sum_{k,l} \underbrace{\left( \frac{\partial \omega_k}{\partial t_l} + (1-\theta_l)\frac{\partial W_l}{\partial \lambda_k} \right)}_{=0} dt_l \wedge d\lambda_k \\
    &\quad + \frac{1}{2}\sum_{m < k} \underbrace{\left( \frac{\partial \omega_k}{\partial \lambda_m} - \frac{\partial \omega_m}{\partial \lambda_k} \right)}_{=0} d\lambda_m \wedge d\lambda_k - \frac{1}{2}\sum_{m < l} \underbrace{\left( (1-\theta_l)\frac{\partial W_l}{\partial t_m} - (1-\theta_m)\frac{\partial W_m}{\partial t_l} \right)}_{=0} dt_m \wedge dt_l \\
    &= 0
\end{align*}
である。

定理3.2が示している多時間ハミルトン系は，自由度 $g$，すなわち $2g$ 個の正準変数 $\lambda = (\lambda_1, \cdots, \lambda_g), \mu = (\mu_1, \cdots, \mu_g)$ の $g$ 次元時間 $t = (t_1, \cdots, t_g)$ に関する微分方程式系である．これら $3g$ 個の変数と $g$ 個のハミルトニアン $H = (H_1, \cdots, H_g$) からなる組
\begin{equation*}
    \mathcal{H} = (\lambda, \mu, H, t)
\end{equation*}
を，$g$-次ガルニエ系のハミルトン構造という．

\end{story}

\begin{Q}[【核心】このページは何を主張しているのか？]
\textbf{A.} 先ほどあなたが「クロスタームが吸収された」と見抜いた、その「吸収された余分な項」が $W_l$ と $\omega_k$ です。
このページは、「変数 $\mu$ を $\nu$ に、$H$ を $K$ にズラすという変換（ゲージ変換の代数的な帰結）は、解析力学における**正準変換（Canonical Transformation）**になっている」という極めて重要な事実を証明しています。
\end{Q}

\begin{Q}[【深掘り】$\wedge$（ウェッジ積）の式は何を意味している？]
\textbf{A.} この $\sum d\mu \wedge d\lambda - \sum dH \wedge dt$ という量は、解析力学において**「ポアンカレ・カルタンの不変量（拡張されたシンプレクティック2形式）」**と呼ばれるものです。
この式が変換の前後で「イコール」で結ばれているということは、「元の系（$H, \mu$）が持っていたハミルトン方程式としての美しい幾何学的構造が、新しい系（$K, \nu$）でも一切壊れずに完全に保存されている」ということを意味します。これが「正準変換」の定義そのものです。
\end{Q}

\begin{Q}[【謎解き】3つ並んだ「偏微分の式」は一体何？]
\textbf{A.} これこそが、正準変換が成り立つための「魔法のタネ明かし」です。
$\nu$ と $K$ をウェッジ積の右辺に代入して左辺と等しくなるためには、「ズレ」である $\omega_k$ や $W_l$ の微分のクロス項が互いに打ち消し合う必要があります。
この3つの式は、「微分すると互いに入れ替わる（渦なし条件、可積分条件）」ことを示しています。これは微分幾何学的に言えば、**「ある1つの巨大なスカラー関数 $S(\lambda, t)$（母関数）が存在して、その勾配（微分）がちょうど $\omega_k$ や $W_l$ になっている」**という事実を保証しています。ズレが「完全微分」で書けるからこそ、シンプレクティック構造が保たれるのです。
\end{Q}

\begin{story}
ハミルトン系に関わる諸概念は，自然なやり方で拡張すれば，$g=1$ の場合がほとんどそのまま通用する．一般に
\begin{equation}
    \Omega = \sum_{k=1}^g d\mu_k \wedge d\lambda_k - \sum_{l=1}^g dH_l \wedge dt_l \label{eq:3.88}
\end{equation}
を多時間ハミルトン系の\textbf{基本2次型式}と呼ぶ．ハミルトン構造の変換
\begin{equation*}
    \mathcal{H} \to \mathcal{H}' = (\lambda', \mu', H', t')
\end{equation*}
が\textbf{正準変換}であるとは，ハミルトン系をハミルトン系に移す，ということである．すなわち
\begin{equation*}
    \Omega' = \sum_{k=1}^g d\mu_k' \wedge d\lambda_k' - \sum_{l=1}^g dH_l' \wedge dt_l'
\end{equation*}
を $\mathcal{H}'$ の基本2次型式とすると，正準変換とは $\Omega = \Omega'$ が成り立つ変換である．

以上の考察により，次の結果が示された．これは後で使われる．
\begin{lemma}[3.7]
変換(3.87)は正準変換である．
\end{lemma}

次に，$a_2^{(l)}(x,t)$ の形を決める．シュレジンガー型の方程式系に関係する命題3.11に対応して，次の補題が成り立つ．
\begin{lemma}[3.8]
$a_2^{(l)}(x,t)$ は次の性質をもつ．
\begin{enumerate}
    \item[(1)] $\xi$ が $t_l$ に依存しないとき，$x=\xi$ で少なくとも1位の零点をもつ．ただし、$x=\infty$の場合は、正則な点である。
    \item[(2)] $x=t_l$ は正則点．
    \item[(3)] $x=\lambda_k$ は高々1位の極。
\end{enumerate}
\end{lemma}

これにより
\begin{equation}
    a_2^{(l)}(x,t) = M_l(t) \frac{T(x)}{L(x)(x-t_l)} \label{eq:3.89}
\end{equation}
とおくことができる．$T(x), L(x)$ は(3.69), (3.70)の多項式である．
\end{story}
実際(1)より、$h\neq l$をみたす$h$について$t_{h}$は$t_{l}$によらないので、$x=t_{h}$で0点を持つから、$\frac{T(x)}{x-t_{l}}$に比例する。(2)から、$x=t_{l}$で極を持たないので、分母の$x-t_{l}$はこれで十分であり、(3)から、$x=\lambda_{k}$は分母には多くても1つしか出てこないことがわかる。よって、分子は$g+2$次式以上で、分母は$g+1$次式以下である。ここで$\xi=\infty$で高々1位の極を持つことに注意すると、(分子)-(分母)は1次以下でなければならない。よって分子が$T(x)$分母が$L(x)(x-t_{l})$であることが確定し、残りは$x$によらない定数倍にしかよらないことがわかる・
\begin{story}
係数 $M_l(t)$ は，方程式(3.82)から決めることができる．すなわち，(3.85), (3.89)を(3.82)に代入して，$(x-t_l)^{-3}$ の係数を比較すると次式が得られる．
\end{story}

\begin{story}
\begin{equation*}
    a_2^{(l)}(t_l, t) = -1
\end{equation*}

したがって，
\begin{equation*}
    M_l(t) = -\frac{L(t_l)}{T'(t_l)} = M_l
\end{equation*}
となる．これは(3.74)で与えられたものと同じである．
\end{story}

\begin{align*}
    r(x,t) &= \frac{\alpha_0}{x^2} + \frac{\alpha_1}{(x-1)^2} + \frac{\alpha_\infty}{x(x-1)} + \sum_{l=1}^g \left[ \frac{\beta_l}{(x-t_l)^2} + \frac{t_l(t_l-1)K_l}{x(x-1)(x-t_l)} \right] \\
    &\quad + \sum_{k=1}^g \left[ \frac{3}{4(x-\lambda_k)^2} - \frac{\lambda_k(\lambda_k-1)\nu_k}{x(x-1)(x-\lambda_k)} \right]
\end{align*}

\begin{equation*}
    a_2^{(l)}(x,t) = M_l(t) \frac{T(x)}{L(x)(x-t_l)} = M_l(t) \frac{x(x-1)\prod_{m \neq l} (x-t_m)}{L(x)}
\end{equation*}

\begin{equation*}
    \frac{\partial^3 a_2^{(l)}}{\partial x^3} - 4r \frac{\partial a_2^{(l)}}{\partial x} - 2\frac{\partial r}{\partial x} a_2^{(l)} + 2\frac{\partial r}{\partial t_l} = 0
\end{equation*}
であるから，最低次の $(x-t_l)^{-3}$ の項を見ると，
\begin{align*}
    & -2 \frac{\partial}{\partial x} \left( \frac{\beta_l}{(x-t_l)^2} \right) \cdot M_l(t) \frac{x(x-1)\prod_{m \neq l} (x-t_m)}{L(x)} + 2 \frac{\partial}{\partial t_l} \left( \frac{\beta_l}{(x-t_l)^2} \right) \\
    &= \frac{\beta_l}{(x-t_l)^3} \left( 4 M_l(t) \frac{x(x-1)\prod_{m \neq l} (x-t_m)}{L(x)} + 4 \right)
\end{align*}

よって，$\displaystyle \lim_{x \to t_l} (x-t_l)^3 \left( \frac{\partial^3 a_2^{(l)}}{\partial x^3} - 4r \frac{\partial a_2^{(l)}}{\partial x} - 2\frac{\partial r}{\partial x} a_2^{(l)} + 2\frac{\partial r}{\partial t_l} \right) = 0$ より，
\begin{equation*}
    \lim_{x \to t_l} 4\beta_l \left( a_2^{(l)}(x,t) + 1 \right) = 0
\end{equation*}

一般に $\beta_l \neq 0$ なので，
\[
\lim_{x \to t_l} a_2^{(l)}(x,t) = -1
\]

よって，
\[
M_l(t) = - \frac{L(t_l)}{T'(t_l)} = M_l
\]
である．

\begin{story}
\begin{proof}[補題3.8の証明]
命題3.8より、$b_{2}^{(l)}=a_{2}^{(l)}$なので、(3.25)について考えればよい．(3.25)の解の基本系 $\vec{z} = \vec{\psi}(x,t)$ で，そのモノドロミー群が $t$ に依らないものをとる．$w(\vec{\psi})$ を，この解のロンスキアンとすると，微分方程式(3.25)の形と、命題2.2から、これは $x$ には依らない．一方，
\[
\frac{\partial \vec{\psi}}{\partial t_{l}} = a_{1}^{(l)}(x,t)\vec{\psi}+a_{2}^{(l)}(x,t)\frac{\partial\vec{\psi}}{\partial x} \quad (l=1,\cdots,L)
\]
から
\begin{equation}
    w(\vec{\psi}) \cdot a_2^{(l)}(x,t) = \det \left( \vec{\psi}(x,t), \frac{\partial}{\partial t_l}\vec{\psi}(x,t) \right) \label{eq:3.90}
\end{equation}
となる．
実際、
\begin{align*}
\det \left( \vec{\psi}, \frac{\partial}{\partial t_l}\vec{\psi}\right)&=\det \left( \vec{\psi}, a_{1}^{(l)}\vec{\psi}+a_{2}^{(l)}\frac{\partial\vec{\psi}}{\partial x}\right)\\
&=\det \left( \vec{\psi}, a_{1}^{(l)}\vec{\psi} \right)+\det \left(\vec{\psi}, a_{2}^{(l)}\frac{\partial\vec{\psi}}{\partial x}\right)\\
&=a_{2}^{(l)}\det \left(\vec{\psi}, \frac{\partial\vec{\psi}}{\partial x}\right)
\end{align*}
$x=x_0$ が(3.25)の正則点ならば，有理関数 $a_2^{(l)}(x,t)$ も $x=x_0$ で正則である．
実際、
\[
a_{2}^{(l)}=\frac{\det \left( \vec{\psi}, \frac{\partial\vec{\psi}}{\partial t_l}\right)}{w(\vec{\psi})}
\]
であり、分母は定数で0でないので極にはなりえず、$\frac{\partial\vec{\psi}}{\partial t_l}$は$x=x_{0}$で正則なので、（コーシーの積分定理等を用いる。）分子も$x=x_{0}$で正則であることからわかる。

特異点 $x=\xi$ ($\infty$を除く)の近傍における $a_2^{(l)}(x,t)$ の様子を局所的に調べる．まず以下の$\psi(x,t)$に対して$\frac{\det \left( \vec{\psi}, \frac{\partial\vec{\psi}}{\partial t_l}\right)}{w(\vec{\psi})}$の値を考える。
\begin{equation*}
    \vec{\psi}(x,t) = \left( (x-\xi)^\sigma(1+\mathcal{O}(x-\xi)), \quad (x-\xi)^{1-\sigma}(1+\mathcal{O}(x-\xi)) \right)
\end{equation*}
をとる．(3.25)の形から特性指数の和が1になることがわかる。この解の基本系 $\vec{\psi}$ の，$x=\xi$ のまわりの局所モノドロミーは $t$ に依らず，また
\begin{equation*}
    w(\vec{\psi}) = 1-2\sigma, \quad \det \left( \vec{\psi}, \frac{\partial \vec{\psi}}{\partial t_l} \right) = (2\sigma - 1)\frac{d\xi}{dt_l} + \mathcal{O}(x-\xi)
\end{equation*}
となる．実際、$w(\vec{\psi})$ は次のように計算できる。
\begin{align*}
    w(\vec{\psi}) &= 
    \begin{vmatrix}
        (x-\xi)^\sigma \big(1 + \mathcal{O}(x-\xi)\big) & (x-\xi)^{1-\sigma} \big(1 + \mathcal{O}(x-\xi)\big) \\[6pt]
        (x-\xi)^{\sigma-1} \big(\sigma + \mathcal{O}(x-\xi)\big) & (x-\xi)^{-\sigma} \big(1-\sigma + \mathcal{O}(x-\xi)\big)
    \end{vmatrix} \\[8pt]
    &= \big(1 - \sigma + \mathcal{O}(x-\xi)\big) - \big(\sigma + \mathcal{O}(x-\xi)\big) \\[4pt]
    &= 1 - 2\sigma + \mathcal{O}(x-\xi)
\end{align*}
ここで、$w(\vec{\psi})$ は定数であることより、$\mathcal{O}(x-\xi)$ の項は消滅し、厳密に以下の値をとる。
\begin{equation}
    w(\vec{\psi}) = 1 - 2\sigma
\end{equation}
また、$\det\left(\vec{\psi}, \frac{\partial \vec{\psi}}{\partial t_l}\right)$について、
\begin{align*}
    \det\left(\vec{\psi}, \frac{\partial \vec{\psi}}{\partial t_l}\right) &= 
    \begin{vmatrix}
        (x-\xi)^\sigma \big(1 + \mathcal{O}(x-\xi)\big) & (x-\xi)^{1-\sigma} \big(1 + \mathcal{O}(x-\xi)\big) \\[6pt]
        (x-\xi)^{\sigma-1} \left(-\sigma \frac{\partial \xi}{\partial t_l} + \mathcal{O}(x-\xi)\right) & (x-\xi)^{-\sigma} \left(-(1-\sigma) \frac{\partial \xi}{\partial t_l} + \mathcal{O}(x-\xi)\right)
    \end{vmatrix} \\[8pt]
    &= -(1-\sigma) \frac{\partial \xi}{\partial t_l} + \sigma \frac{\partial \xi}{\partial t_l} + \mathcal{O}(x-\xi) \\[4pt]
    &= (2\sigma - 1) \frac{\partial \xi}{\partial t_l} + \mathcal{O}(x-\xi)
\end{align*}

ただし、この計算において「モノドロミーは $t$ に依存しないので、特性指数 $\sigma$ は $t$ に依存しない」という事実を用いている。
仮定(P2)から $1-2\sigma \neq 0$ であり，（実際$\sigma=\frac{1}{2}$なら、特性指数がともに$\frac{1}{2}$となって差は0）したがって、$x=\xi$ の近傍で
\begin{equation*}
    \frac{\det \left( \vec{\psi}, \frac{\partial\vec{\psi}}{\partial t_l}\right)}{w(\vec{\psi})}=-\frac{\partial \xi}{\partial t_l} + \mathcal{O}(x-\xi)
\end{equation*}
と表される。

$x=\infty$ の局所座標を $x = 1/u$ とする。
\begin{equation}
    - \frac{\det\left(\vec{\psi}, \frac{\partial \vec{\psi}}{\partial t_l}\right)}{\det\left(\vec{\psi}, u^2 \frac{\partial \vec{\psi}}{\partial u}\right)} \label{eq:a2_inf_2}
\end{equation}
の値を求める。$u=0$ のまわりでの基本解行列 $\vec{\psi}(u,t)$ の漸近展開を次のように置く。
\begin{align*}
    \vec{\psi}_1(u,t) &= u^\sigma \big(1 + c(t)u + \mathcal{O}(u^2)\big) \\
    \vec{\psi}_2(u,t) &= u^{-1-\sigma} \big(1 + d(t)u + \mathcal{O}(u^2)\big)
\end{align*}

分母の行列式（ロンスキアンに対応）は定数となり、以下の通り計算される。
\begin{align*}
    \det\left(\vec{\psi}, u^2 \frac{\partial \vec{\psi}}{\partial u}\right) &= 
    \begin{vmatrix}
        u^\sigma (1 + \mathcal{O}(u)) & u^{-1-\sigma} (1 + \mathcal{O}(u)) \\[6pt]
        u^{\sigma+1} (\sigma + \mathcal{O}(u)) & u^{-\sigma} (-(1+\sigma) + \mathcal{O}(u))
    \end{vmatrix} \\[8pt]
    &= -(1+\sigma) - \sigma + \mathcal{O}(u) \\
    &= -(1+2\sigma) + \mathcal{O}(u)
\end{align*}

一方、分子の行列式については、時間微分 $\partial_{t_l}$ が $u^\sigma$ 等の主要部を消去することに注意すると、
\begin{align*}
    \det\left(\vec{\psi}, \frac{\partial \vec{\psi}}{\partial t_l}\right) &= 
    \begin{vmatrix}
        u^\sigma \big(1 + c(t)u + \mathcal{O}(u^2)\big) & u^{-1-\sigma} \big(1 + d(t)u + \mathcal{O}(u^2)\big) \\[6pt]
        u^\sigma \left(\frac{\partial c(t)}{\partial t_l}u + \mathcal{O}(u^2)\right) & u^{-1-\sigma} \left(\frac{\partial d(t)}{\partial t_l}u + \mathcal{O}(u^2)\right)
    \end{vmatrix} \\[8pt]
    &= u^{-1} \left[ \big(1 + c(t)u\big)\left(\frac{\partial d(t)}{\partial t_l}u\right) - \big(1 + d(t)u\big)\left(\frac{\partial c(t)}{\partial t_l}u\right) \right] + \mathcal{O}(u) \\[4pt]
    &= \frac{\partial d(t)}{\partial t_l} - \frac{\partial c(t)}{\partial t_l} + \mathcal{O}(u)
\end{align*}

よって、式\eqref{eq:a2_inf_2}に代入すると、
\begin{equation}
     - \frac{\det\left(\vec{\psi}, \frac{\partial \vec{\psi}}{\partial t_l}\right)}{\det\left(\vec{\psi}, u^2 \frac{\partial \vec{\psi}}{\partial u}\right)}= \frac{\frac{\partial d(t)}{\partial t_l} - \frac{\partial c(t)}{\partial t_l}}{1+2\sigma} + \mathcal{O}(u)
\end{equation}
となる。

大域的な基本解系を $\vec{Y}(x,t)$、特異点 $x=\xi$ における局所正規化解を $\vec{\psi}(x,t) = (\psi_1, \psi_2)$ とし、接続行列を $C(t) = (C_{ij}(t))$ として $\vec{Y} = \vec{\psi}C(t)$ と表す。

$a_2^{(l)}$ の計算に必要な行列式を展開する。まず分母のロンスキアンは、
\begin{equation}
    W(\vec{Y}) = \det(\vec{\psi}C, \partial_x(\vec{\psi}C)) = \det C(t) W(\vec{\psi})
\end{equation}
分子の時間微分に関する行列式は、分配則により次のように展開される。
\begin{equation}
    \det(\vec{Y}, \partial_{t_l}\vec{Y}) = \det C \det(\vec{\psi}, \partial_{t_l}\vec{\psi}) + \det(\vec{\psi}C, \vec{\psi}\partial_{t_l}C)
\end{equation}
右辺第2項を成分で具体的に計算すると、
\begin{align*}
    \det(\vec{\psi}C, \vec{\psi}\partial_{t_l}C) 
    &= \begin{vmatrix} 
        C_{11}\psi_1 + C_{21}\psi_2 & C_{12}\psi_1 + C_{22}\psi_2 \\ 
        \dot{C}_{11}\psi_1 + \dot{C}_{21}\psi_2 & \dot{C}_{12}\psi_1 + \dot{C}_{22}\psi_2 
    \end{vmatrix} \\
    &= A\psi_1^2 + B\psi_2^2 + D\psi_1\psi_2
\end{align*}
ここで各係数は以下の通りである。
\begin{align*}
    A &= \begin{vmatrix} C_{11} & C_{12} \\ \dot{C}_{11} & \dot{C}_{12} \end{vmatrix} = C_{11}\dot{C}_{12} - \dot{C}_{11}C_{12} \\
    B &= \begin{vmatrix} C_{21} & C_{22} \\ \dot{C}_{21} & \dot{C}_{22} \end{vmatrix} = C_{21}\dot{C}_{22} - \dot{C}_{21}C_{22} \\
    D &= \begin{vmatrix} C_{21} & C_{12} \\ \dot{C}_{21} & \dot{C}_{12} \end{vmatrix} + \begin{vmatrix} C_{11} & C_{22} \\ \dot{C}_{11} & \dot{C}_{22} \end{vmatrix}
\end{align*}

局所解 $\vec{\psi}$ の $x=\xi$ 周りのモノドロミー行列を $\Gamma = \mathrm{diag}(e^{2\pi i \sigma}, e^{-2\pi i \sigma})$ とする。大域解 $\vec{Y}$ のモノドロミー行列は $C^{-1}\Gamma C$ となる。
等モノドロミー変形の要請により、この行列は時間 $t_l$ に依存してはならない。
\begin{equation}
    C^{-1}\Gamma C = \frac{1}{\det C} 
    \begin{pmatrix} 
        C_{11}C_{22}e^{2\pi i \sigma} - C_{12}C_{21}e^{-2\pi i \sigma} & C_{12}C_{22}(e^{2\pi i \sigma} - e^{-2\pi i \sigma}) \\ 
        -C_{11}C_{21}(e^{2\pi i \sigma} - e^{-2\pi i \sigma}) & C_{11}C_{22}e^{-2\pi i \sigma} - C_{12}C_{21}e^{2\pi i \sigma} 
    \end{pmatrix}
\end{equation}
非対角成分が時間に依存しないことから、以下の微分方程式が成立する。
\begin{equation}
    \frac{d}{dt_l} \left( \frac{C_{12}C_{22}}{\det C} \right) = 0, \quad \frac{d}{dt_l} \left( \frac{C_{11}C_{21}}{\det C} \right) = 0
\end{equation}
第一式より、
\begin{equation*}
    \frac{d}{dt_l} \left( \frac{1}{\frac{C_{11}}{C_{12}} - \frac{C_{21}}{C_{22}}} \right) = -\frac{1}{\left(\frac{C_{11}}{C_{12}} - \frac{C_{21}}{C_{22}}\right)^2} \left( \frac{d}{dt_l}\left(\frac{C_{11}}{C_{12}}\right) - \frac{d}{dt_l}\left(\frac{C_{21}}{C_{22}}\right) \right) = 0
\end{equation*}
ここで $\frac{d}{dt_l}\left(\frac{C_{11}}{C_{12}}\right) = \frac{-A}{C_{12}^2}$、$\frac{d}{dt_l}\left(\frac{C_{21}}{C_{22}}\right) = \frac{-B}{C_{22}^2}$ であるから、
\begin{equation}
    \frac{A}{C_{12}^2} = \frac{B}{C_{22}^2} \label{eq:mono1}
\end{equation}
同様に、第二式から商の微分を用いて整理すると、
\begin{equation}
    \frac{A}{C_{11}^2} = \frac{B}{C_{21}^2} \label{eq:mono2}
\end{equation}
式\eqref{eq:mono1}および式\eqref{eq:mono2}から $A$ を消去すると、
\begin{equation}
    \left( \frac{C_{11}^2}{C_{21}^2} - \frac{C_{12}^2}{C_{22}^2} \right) B = \left( \frac{C_{11}}{C_{21}} + \frac{C_{12}}{C_{22}} \right) \left( \frac{C_{11}}{C_{21}} - \frac{C_{12}}{C_{22}} \right) B = 0
\end{equation}
$\det C \neq 0$ より $\frac{C_{11}}{C_{21}} - \frac{C_{12}}{C_{22}} \neq 0$ である。非自明な接続行列においては第一因数も一般にゼロにならないため、$B=0$ を得る。背理的に $A=0$ も従う。
ここで、第一因数が $0$ となるエッジケース、すなわち $\frac{C_{11}}{C_{21}} = - \frac{C_{12}}{C_{22}}$ のときを考える。このとき $\frac{C_{11}}{C_{12}} = -\frac{C_{21}}{C_{22}}$ となるため、式\eqref{eq:mono1}の導出元である微分方程式に代入すると、
\begin{equation*}
    2 \frac{d}{dt_l}\left(\frac{C_{11}}{C_{12}}\right) = 0 \quad \Longrightarrow \quad \frac{C_{11}'C_{12} - C_{11}C_{12}'}{C_{12}^2} = 0
\end{equation*}
ここで、$C_{12} = 0 \Rightarrow A=0$ より $B=0$ で題意を満たす。
また $C_{12} \neq 0$ ならば $A=0$ となり、式\eqref{eq:mono1}より結果として $B=0$ を得る。
以上より、いかなる場合においても $A=0, B=0$ が成立し、多価の項は完全に消滅する。
以上により、有理関数を阻害する多価の項 $\psi_1^2$ と $\psi_2^2$ の係数は恒等的に消滅することが代数的に証明された。

$A=B=0$ の結果を用いると、$a_2^{(l)}$ の厳密な表示は次のように簡約化される。
\begin{equation}
    a_2^{(l)} = \frac{\det C \det(\vec{\psi}, \partial_{t_l}\vec{\psi}) + D\psi_1\psi_2}{\det C \cdot W(\vec{\psi})} = \frac{\det(\vec{\psi}, \partial_{t_l}\vec{\psi})}{W(\vec{\psi})} + \frac{D\psi_1\psi_2}{\det C \cdot W(\vec{\psi})}
\end{equation}

\textbf{(1) 固定特異点 $x=\xi$ の場合} \\
$\xi$ は固定されているため $\partial_{t_l}\vec{\psi}$ の主要部は消滅し、第一項は $\mathcal{O}((x-\xi)^2)$ となる。一方、第二項は $\psi_1\psi_2 \sim (x-\xi)^\sigma (x-\xi)^{1-\sigma} = (x-\xi)^1$ である。
したがって、主要部は第二項に支配され、
\begin{equation}
    a_2^{(l)} \sim \frac{D}{\det C \cdot W(\vec{\psi})} (x-\xi) = \mathcal{O}(x-\xi)
\end{equation}
となり、厳密に1位の零点を持つことが示された。

\textbf{(2) 無限遠点 $x=\infty$ の場合} \\
局所座標 $x=1/u$ を導入する。微分演算子 $\partial_x = -u^2 \partial_u$ の寄与により、分母に $u^2$ の補正が生じる。$\psi_1\psi_2 \sim u^{-\sigma} u^{\sigma+1} = u$ であるため、
\begin{align*}
    a_2^{(l)} &= \frac{\det\left(C\vec{\psi}, \frac{\partial (C\vec{\psi})}{\partial t_l}\right)}{\det\left(C\vec{\psi}, u^2 \frac{\partial (C\vec{\psi})}{\partial u}\right)} \\[8pt]
    &= \frac{\det\left(C\vec{\psi}, C\frac{\partial \vec{\psi}}{\partial t_l}\right)}{\underbrace{\det\left(\vec{\psi}, u^2 \frac{\partial \vec{\psi}}{\partial u}\right)}_{\mathcal{O}(1)}} + \frac{1}{u^2} \frac{\overbrace{D\psi_1\psi_2}^{\sim \mathcal{O}(u)}}{\det C \cdot W(\vec{\psi})} \\[8pt]
    &= \mathcal{O}(u^{-1})
\end{align*}
となり、厳密に1位の極を持つことが示された。命題3.8の証明から，これが $a_2^{(l)}(x,t)$ を局所的に定める。これで補題3.8の(1), (2)が示された．

見かけの特異点 $x=\lambda$ の近傍では，基本解を
\begin{equation*}
    \vec{\psi}(x,t) = \left( (x-\lambda)^\sigma(1+\mathcal{O}(x-\lambda)), \quad (x-\lambda)^{1-\sigma}(1+\mathcal{O}(x-\lambda)) \right)
\end{equation*}
と定めると、特性指数は$-\frac{1}{2},\frac{3}{2}$より、$\sigma=\frac{3}{2}$として、モノドロミー行列$\Gamma$は$-I$となる。よって変換行列$C(t)$をかけても、モノドロミー行列は$C^{-1}\Gamma C=\Gamma$よりCに課される条件はない。よって動く特異点としての主要部と合わせて、
\begin{align*}
\det(\vec{\psi}C, \vec{\psi}\partial_{t_l}C) &= A\psi_1^2 + B\psi_2^2 + D\psi_1\psi_2\\
&=A\mathcal{O}((x-\lambda)^{3}) + B\mathcal{O}((x-\lambda)^{-1}) + D\mathcal{O}(x-\lambda)
\end{align*}
となることがわかる。よって、
\begin{align*}
a_{2}^{(l)}(x,t)&=\frac{\partial \lambda}{\partial t_l}+\mathcal{O}((x-\lambda)^{-1})\\
&=\mathcal{O}((x-\lambda)^{-1})
\end{align*}
このことから補題3.8の(3)が従う．
\end{proof}
\end{story}
\begin{proposition}[3.8]
非斉次微分方程式(3.29)を満たす，$x$ の有理関数は存在するとしてもただ1つである
\end{proposition}

\begin{proposition}[2.2]
    \[
    \frac{d^{2}y}{dx^{2}}+p_{1}(x)\frac{dy}{dx}+p_{2}(x)=0
    \]
    の解の基本系$\vec{\varphi}(x)$のロンスキアンは
    \[
    \frac{dw}{dx}+p_{1}(x)w=0
    \]
    なる$1$階微分方程式を満たす。
\end{proposition}

\begin{story}

最後に，(3.85), (3.89) を，微分方程式系(3.82), (3.83) に代入することにより，次の補題が証明される．

% --- 補題3.9 ---
\begin{lemma}[3.9]
    (3.85), (3.89) が成立するための条件は，アクセサリー・パラメータ $\lambda = (\lambda_1, \cdots, \lambda_g)$, $\nu = (\nu_1, \cdots, \nu_g)$ が，$K = (K_1, \cdots, K_g)$ をハミルトニアンとするハミルトン系
    \begin{equation}
        \frac{\partial \lambda_k}{\partial t_l} = \frac{\partial K_l}{\partial \nu_k}, \quad 
        \frac{\partial \nu_k}{\partial t_l} = - \frac{\partial K_l}{\partial \lambda_k} \quad 
        (k,l = 1, \cdots, g)
        \tag{3.91}\label{eq:3.91}
    \end{equation}
    を満たすことである．
\end{lemma}

証明の方針は、係数の比較だけだが、計算はめんどうで若干の技巧を要する。

微分方程式系\eqref{eq:3.91}は、(3.82)だけから導かれる。(3.83)からは
\begin{equation}
    \frac{\lambda_k - t_l}{M_l} \frac{\partial \lambda_k}{\partial t_l} - \frac{\lambda_k - t_h}{M_h} \frac{\partial \lambda_k}{\partial t_h} 
    = \frac{(t_l - t_h) T(\lambda_k)}{(\lambda_k - t_l)(\lambda_k - t_h) L'(\lambda_k)} \quad 
    (h, k, l = 1, \cdots, g)
    \label{eq:3.92}
\end{equation}
という式が出るが、命題3.14に示した通り、これは(3.82)の積分可能条件であり、この意味で\eqref{eq:3.91}に含まれる。第3.8節166ページ、(3.109)において、\eqref{eq:3.92}の意味について述べる。

計算の見通しを良くするために、(3.83)から\eqref{eq:3.92}を直接導いてみよう。

\begin{equation}
    A_{hl} = \frac{1}{a_2^{(h)} a_2^{(l)}} \left( \frac{\partial a_2^{(l)}}{\partial t_h} - \frac{\partial a_2^{(h)}}{\partial t_l} \right) - \frac{\partial}{\partial x} \log \frac{a_2^{(l)}}{a_2^{(h)}} \notag
\end{equation}
とおく。(3.83) は $A_{hl} = 0$ と同じである。

\end{story}
実際、ベクトル場の可換性（積分可能条件）より、
\begin{align*}
    \left(\frac{\partial}{\partial t_h} - a_2^{(h)} \frac{\partial}{\partial x}\right) a_2^{(l)} &= \left(\frac{\partial}{\partial t_l} - a_2^{(l)} \frac{\partial}{\partial x}\right) a_2^{(h)} \\
    \frac{\partial a_2^{(l)}}{\partial t_h} - \frac{\partial a_2^{(h)}}{\partial t_l} - a_2^{(h)} \frac{\partial a_2^{(l)}}{\partial x} + a_2^{(l)} \frac{\partial a_2^{(h)}}{\partial x} &= 0
\end{align*}
両辺を $a_2^{(l)} a_2^{(h)}$ で割って整理すると、
\begin{equation*}
    \frac{1}{a_2^{(l)}a_2^{(h)}} \left( \frac{\partial a_2^{(l)}}{\partial t_h} - \frac{\partial a_2^{(h)}}{\partial t_l} \right) - \frac{1}{a_2^{(l)}} \frac{\partial a_2^{(l)}}{\partial x} + \frac{1}{a_2^{(h)}} \frac{\partial a_2^{(h)}}{\partial x} = 0
\end{equation*}
すなわち、$A_{hl}$ の定義式と一致し、
\begin{equation}
    A_{hl} = \frac{1}{a_2^{(h)}} \frac{\partial}{\partial t_h} \log a_2^{(l)} - \frac{1}{a_2^{(l)}} \frac{\partial}{\partial t_l} \log a_2^{(h)} - \frac{\partial}{\partial x} \log \frac{a_2^{(l)}}{a_2^{(h)}} = 0 \label{eq:A_hl_def}
\end{equation}
を得る。

\begin{story}

\begin{lemma}
    $A_{hl}$ は $x$ の有理関数で，$g$ 個の点 $x=0, 1, t_p$ $(p \neq h, l)$ で高々1位の極，$x=\infty$ で少なくとも1位の零点をもつ．
\end{lemma}

\begin{proof}
    補題3.8により
    \begin{equation*}
        a_2^{(l)}(x,t) = M_l \frac{T(x)}{L(x)(x-t_l)}, \quad M_l = -\frac{L(t_l)}{T'(t_l)}
    \end{equation*}
    であったから，$A_{hl}$ の定義式に含まれる各対数微分を計算する。
\begin{align*}
    \frac{\partial}{\partial x} \log a_2^{(l)} &= \frac{\partial}{\partial x} \left( \log \left( M_l \frac{T(x)}{L(x)(x-t_l)} \right) \right) \\
    &= \frac{\partial}{\partial x} \left( \log M_l + \log x + \log(x-1) + \sum_{i \neq l} \log(x-t_i) - \sum_{k=1}^g \log(x-\lambda_k) \right) \\
    &= \frac{1}{x} + \frac{1}{x-1} + \sum_{i \neq l} \frac{1}{x-t_i} - \sum_{k=1}^g \frac{1}{x-\lambda_k} \\[8pt]
    \frac{\partial}{\partial t_h} \log a_2^{(l)} &= \frac{\partial}{\partial t_h} \left( \log \left( M_l \frac{T(x)}{L(x)(x-t_l)} \right) \right) \\
    &= \frac{\partial}{\partial t_h} \left( \log M_l + \log x + \log(x-1) + \sum_{i \neq l} \log(x-t_i) - \sum_{k=1}^g \log(x-\lambda_k) \right) \\
    &= \frac{\partial}{\partial t_h} \log M_l - \frac{1}{x-t_h} + \sum_{k=1}^g \frac{1}{x-\lambda_k} \frac{\partial \lambda_k}{\partial t_h}
\end{align*}

これを用いて $A_{hl}$ の第一項に代入する。
\begin{align*}
    \frac{1}{a_2^{(h)}} \frac{\partial}{\partial t_h} \log a_2^{(l)} 
    &= \frac{(x-t_h)L(x)}{M_h T(x)} \left( \frac{\partial}{\partial t_h} \log M_l - \frac{1}{x-t_h} + \sum_{k=1}^g \frac{1}{x-\lambda_k} \frac{\partial \lambda_k}{\partial t_h} \right) \\
    &= \frac{(x-t_h)L(x)}{M_h T(x)} \frac{\partial}{\partial t_h} \log M_l - \frac{1}{M_h} \frac{L(x)}{T(x)} + \frac{1}{M_h} \sum_{k=1}^g \frac{x-t_h}{T(x)} \frac{L(x)}{x-\lambda_k} \frac{\partial \lambda_k}{\partial t_h}
\end{align*}

これを $A_{hl}$ の定義式全体に代入して整理すると、
\begin{align*}
    A_{hl} &= \frac{L(x)}{T(x)} \left\{ \frac{x-t_h}{M_h} \frac{\partial}{\partial t_h} \log M_l - \frac{x-t_l}{M_l} \frac{\partial}{\partial t_l} \log M_h - \frac{1}{M_h} + \frac{1}{M_l} \right. \\
    &\quad \left. + \sum_{k=1}^g \frac{1}{x-\lambda_k} \left( \frac{x-t_h}{M_h} \frac{\partial \lambda_k}{\partial t_h} - \frac{x-t_l}{M_l} \frac{\partial \lambda_k}{\partial t_l} \right) \right\} + \frac{1}{x-t_l} - \frac{1}{x-t_h}
\end{align*}
よって、有理関数であり、係数 $1/T(x)$ の寄与から、$x=0, 1, t_p\ (p \neq l,h)$ は高々1位の極をもつ。

また、無限遠点での挙動を見るため $x=\frac{1}{u}$ を代入する。
\begin{align*}
    A_{hl} &= u \frac{\prod (1-\lambda_k u)}{(1-u) \prod (1-t_m u)} \left\{ (1-t_h u) \frac{\partial}{\partial t_h} \log M_l - (1-t_l u) \frac{\partial}{\partial t_l} \log M_h - \frac{u}{M_h} + \frac{u}{M_l} \right. \\
    &\quad \left. + u \sum_{k=1}^g \frac{1}{1-\lambda_k u} \left( (1-t_h u) \frac{\partial \lambda_k}{\partial t_h} - (1-t_l u) \frac{\partial \lambda_k}{\partial t_l} \right) \right\} + \frac{u}{1-t_l u} - \frac{u}{1-t_h u} \\
    &=: u B_{hl}
\end{align*}
ここで $u \to 0$ の極限をとると、
\begin{equation*}
    B_{hl} \xrightarrow{u \to 0} \frac{\partial}{\partial t_h} \log M_l - \frac{\partial}{\partial t_l} \log M_h
\end{equation*}
となるため、$A_{hl}$ は $x=\infty$ で少なくとも1位の零点をもつ。
\end{proof}
\end{story}

ここで、$x=t_{h},t_{l}ではA_{hl}$は正則である。実際、$x=t_h$ 近傍における $A_{hl}$ の主要部を評価し、留数 (Residue) を計算すると
$A_{hl}$ の各項のうち、$x=t_h$ で極を持ちうるのは以下の部分である。
\begin{equation*}
    \frac{1}{a_2^{(h)}} \frac{\partial}{\partial t_h} \log a_2^{(l)} \sim - \frac{1}{a_2^{(h)}} \frac{1}{x-t_h}, \qquad - \frac{\partial}{\partial x} \log a_2^{(l)} \sim - \frac{1}{x-t_h}
\end{equation*}
ここで、$a_2^{(h)}$ は $x=t_h$ を分母の $L(x)$ 因子に持つため、テキストの条件より $x \to t_h$ として極限（留数）が $a_2^{(h)}(t_h, t) \to -1$ となる。
したがって、$x=t_h$ における留数は次のように計算される。
\begin{equation*}
    \mathop{\mathrm{Res}}_{x=t_h} A_{hl} = - \frac{1}{-1} - 1 = 0
\end{equation*}
留数が $0$ となることから、一見存在するように見えた $x=t_h$ での極は完全に相殺し、正則であることが証明された。$x=t_l$ についても対称性より同様である。

\begin{story}
一方，$A_{hl}(x)$ は $x=\lambda_k$ で正則で
$A_{hl}$ の $x-\lambda_k$ の $0$ 次の項は、
\begin{equation*}
    \frac{L(x)}{x-\lambda_k} \cdot \frac{1}{T(x)} \left( \frac{x-t_h}{M_h} \frac{\partial \lambda_k}{\partial t_h} - \frac{x-t_l}{M_l} \frac{\partial \lambda_k}{\partial t_l} \right) + \frac{1}{x-t_l} - \frac{1}{x-t_h}
\end{equation*}
よって、$x=\lambda_k$ を代入して
\begin{align*}
    A_{hl}(\lambda_k) &= \frac{\prod_{m \neq k} (\lambda_k - \lambda_m)}{\lambda_k(\lambda_k-1)\prod_{m \neq h} (\lambda_k - t_m)} \cdot \left( \frac{1}{M_h} \frac{\partial \lambda_k}{\partial t_h} \right) \\
    &\quad - \frac{\prod_{m \neq k} (\lambda_k - \lambda_m)}{\lambda_k(\lambda_k-1)\prod_{m \neq l} (\lambda_k - t_m)} \cdot \left( \frac{1}{M_l} \frac{\partial \lambda_k}{\partial t_l} \right) + \frac{1}{\lambda_k - t_l} - \frac{1}{\lambda_k - t_h} \\[8pt]
    &= \frac{L'(\lambda_k)}{T(\lambda_k)} \cdot \left( \frac{\lambda_k-t_h}{M_h} \cdot \frac{\partial \lambda_k}{\partial t_h} - \frac{\lambda_k-t_l}{M_l} \frac{\partial \lambda_k}{\partial t_l} \right)
\end{align*}
となる。補題3.10により，零点と極の個数を数えればわかるように，有理関数として $A_{hl}(x) = 0$ となるための必要十分条件は
\begin{equation*}
    A_{hl}(\lambda_k) = 0 \quad (k = 1, 2, \cdots, g)
\end{equation*}
であり，この式を整理して(3.92)が成り立つことがわかった．
\end{story}

実際同値である。

$(\Leftarrow)$ \quad $A_{hl}(x) \equiv 0$ ならば、当然 $A_{hl}(\lambda_k) = 0$ であるから成立する

$(\Rightarrow)$ \quad $A_{hl}(\lambda_k) = 0 \ (k=1,\dots,g)$ とすると、$x=\infty$ での零点と合わせて、少なくとも $g+1$ 個の零点をもつ。また、$A_{hl}(x)$ は $x=0, 1, t_p \ (p \neq l, h)$ では高々1位の極をもち、それ以外の点では正則である。（すなわち、極の総数は高々 $g$ 個である。）ここで、もし $A_{hl}(x) \not\equiv 0$ と仮定する。リーマン球面 $\mathbb{P}^1(\mathbb{C})$ 上の偏角の原理から、0でない有理型関数の零点の総数と極の総数は等しくならなければならない。しかし、$A_{hl}(x)$ の零点は $g+1$ 個以上、極は $g$ 個以下であり、数が一致せず矛盾する。
したがって、$A_{hl}(x) \equiv 0$ である。

\begin{story}
この節の残りで補題3.9の証明の概略を紹介する．本体の条件を解析するために
\begin{equation}
    A_l = \frac{1}{2} a_2^{(l)} \frac{\partial^3 a_2^{(l)}}{\partial x^3} - \frac{\partial}{\partial x} \left( r(x,t) (a_2^{(l)})^2 \right) + a_2^{(l)} \frac{\partial r}{\partial t_l} (x,t) \notag
\end{equation}
とおき，すぐ上の議論と同様に $A_l = 0$ となる条件を絞り込んでいく．

\begin{lemma}
    有理関数 $A_l$ は次の性質をもつ．
    \begin{itemize}
        \item[(1)] $x=0, 1, t_h$ $(h \neq l)$ で正則
        \item[(2)] $x=\lambda_k$ は高々4位の極
        \item[(3)] $x=t_l$ は高々1位の極
        \item[(4)] $x=\infty$ は少なくとも2位の零点
    \end{itemize}
\end{lemma}

\begin{proof}
(3.85)より、$r$はどの確定特異点でも高々2位の極を持つので、$u=x-\xi$ を特異点 $x=\xi$ における局所座標として
\[
    r(x,t) = \frac{\alpha}{u^2} + \frac{\beta}{u} + \cdots
\]
と展開する．$\alpha$は特性指数の積など、$t$によらないものの合成で書けるので、$\frac{\partial \alpha}{\partial t_l}=0$ であるから
\[
    \frac{\partial r}{\partial t_l}(x,t) = \frac{2\alpha}{u^3}\frac{\partial \xi}{\partial t_l} + \frac{\beta}{u^2}\frac{\partial \xi}{\partial t_l} + \frac{1}{u}\frac{\partial \beta}{\partial t_l}+\cdots
\]
もし $\frac{\partial \xi}{\partial t_l}=0$ ならば，このような $x=\xi$ については $a_2^{(l)}(\xi,t)=0$ であるから，
$a_2^{(l)}$ は $x = 0, 1, t_h \ (h \neq l)$ で $1$位の零点をもち、
\begin{align*}
    a_2^{(l)} \frac{\partial^3 a_2^{(l)}}{\partial x^3} &\sim \mathcal{O}(u) \mathcal{O}(1) \sim \mathcal{O}(u) \\[8pt]
    \frac{\partial}{\partial x} \left( r (a_2^{(l)})^2 \right) &\sim \frac{\partial}{\partial x} \left( \mathcal{O}(u^{-2}) \mathcal{O}(u^2) \right) \sim \frac{\partial}{\partial x} (\mathcal{O}(1)) \sim \mathcal{O}(1) \\[8pt]
    a_2^{(l)} \frac{\partial r}{\partial t_l} &\sim \mathcal{O}(u) \mathcal{O}(u^{-1}) \sim \mathcal{O}(1)
\end{align*}
より、
\begin{equation*}
    A_l \sim \mathcal{O}(1)
\end{equation*}
つまり、(1)が従う．

$\xi=t_l$ ならば，$a_2^{(l)}(t_l,t)=-1+\eta u+\cdots$ という形だから
\begin{align*}
    a_2^{(l)} \frac{\partial^3 a_2^{(l)}}{\partial x^3} &\sim \mathcal{O}(1) \mathcal{O}(1) \sim \mathcal{O}(1) \\
    a_2^{(l)} \frac{\partial r}{\partial t_l}(x,t) &= -\frac{2\alpha}{u^3} + \frac{2\alpha\eta-\beta}{u^2} + \mathcal{O}(u^{-1}) \\
    r(x,t)\left(a_2^{(l)}\right)^2 &= \frac{\alpha}{u^2} + \frac{\beta-2\alpha\eta}{u} + \mathcal{O}(1)
\end{align*}
となる．ここで $\mathcal{O}(1)$ は $u=0$ で正則な部分を表す．これを代入すれば，$x=t_l$ の近傍で、高次の極が打ち消しあって、
\[
    A_l = \mathcal{O}(u^{-1})
\]
となり，(3)が示された．

次に $u=\frac{1}{x}$ とすると，$u=0$ の近傍で $\frac{\partial}{\partial x}=-u^2 \frac{\partial}{\partial u}$ より
\begin{align*}
    \frac{1}{M_l}a_2^{(l)} &= \frac{1}{u}+a+bu+\cdots, & \frac{1}{M_l}\frac{\partial^3 a_2^{(l)}}{\partial x^3} &= -6bu^4+\cdots \\
    r(x,t) &= cu^2+fu^3+\cdots, & \frac{\partial r}{\partial t_l}(x,t) &= \frac{\partial f}{\partial t_l}u^3, \quad \frac{\partial c}{\partial t_l}=0
\end{align*}
と表される。実際

\begin{align*}
\frac{1}{M_l} a_2^{(l)} &= \frac{x(x-1) \prod_{k \neq l} (x-t_k)}{\prod_{k=1}^g (x-\lambda_k)} \\
&\quad x = \frac{1}{u} \text{として、} \\
\frac{1}{M_l} a_2^{(l)} &= \frac{1}{u} \underbrace{\frac{(1-u) \prod_{k \neq l} (1-t_k u)}{\prod_{k=1}^g (1-\lambda_k u)}}_{=: B_l}
\end{align*}

$B_l \xrightarrow{u \to 0} 1$ より、
\begin{equation*}
\frac{1}{M_l} a_2^{(l)} = \frac{1}{u} + a + b u + \mathcal{O}(u^2)
\end{equation*}

また、$\frac{\partial}{\partial x} = -u^2 \frac{\partial}{\partial u}$ より、
\begin{align*}
\frac{1}{M_l} \frac{\partial^3 a_2^{(l)}}{\partial x^3} = \frac{\partial^3}{\partial x^3} \left( \frac{a_2^{(l)}}{M_l} \right) &= - \left( u^2 \frac{\partial}{\partial u} \right)^3 \left( \frac{1}{u} + a + b u + \mathcal{O}(u^2) \right) \\
&= - \left( u^2 \frac{\partial}{\partial u} \right)^2 \left( -1 + b u^2 + \mathcal{O}(u^3) \right) \\
&= - \left( u^2 \frac{\partial}{\partial u} \right) \left( 2 b u^3 + \mathcal{O}(u^4) \right) \\
&= -6 b u^4 + \mathcal{O}(u^5)
\end{align*}

また、
\begin{align*}
r\left(\frac{1}{u}, t\right) = u^2& \left( \alpha_0 + \frac{\alpha_1}{(1-u)^2} + \frac{\alpha_\infty}{1-u} + \sum_{i=1}^g \left[ \frac{\beta_i}{(1-t_i u)^2} + \frac{u t_i(t_i-1) K_i}{(1-u)(1-t_i u)} \right]\right. \\
&\left.+ \sum_{k=1}^g \left[ \frac{\gamma_k}{(1-\lambda_k u)^2} - \frac{u \lambda_k(\lambda_k-1) \nu_k}{(1-u)(1-\lambda_k u)} \right] \right) 
\end{align*}
より$r$ は $x = \infty$ を $2$位の零点にもつ。よって、$r\left(\frac{1}{u}, t\right) = c u^2 + f u^3 + \mathcal{O}(u^4)$ここで $c = \alpha_0 + \alpha_1 + \alpha_\infty + \sum_{i=1}^g \beta_i + \sum_{k=1}^g \gamma_k = \frac{1-\kappa_\infty^2}{4}$ より、$\kappa_\infty$ は SL型に変形する前の特性指数で $t$ によらない。
つまり $\frac{\partial c}{\partial t_l} = 0$ ．よって、
\begin{equation*}
\frac{\partial r}{\partial t_l} = \frac{\partial f}{\partial t_l} u^3 + \mathcal{O}(u^4)
\end{equation*}

以上より、
\begin{align*}
a_2^{(l)} \frac{\partial^3 a_2^{(l)}}{\partial x^3} &\sim \mathcal{O}(u^{-1}) \mathcal{O}(u^4) \sim \mathcal{O}(u^3) \\
a_2^{(l)} \frac{\partial r}{\partial t_l} &\sim \mathcal{O}(u^{-1}) \mathcal{O}(u^3) \sim \mathcal{O}(u^2) \\
\frac{\partial}{\partial x} \left( r (a_2^{(l)})^2 \right) &\sim -u^2 \frac{\partial}{\partial u} \left( \mathcal{O}(u^2) \mathcal{O}(u^{-2}) \right) \sim \mathcal{O}(u^2)
\end{align*}

よって、 $A_l \sim \mathcal{O}(u^2)$、つまり(4)がわかる．

最後に $u=x-\lambda_k$ として
\begin{equation*}
    a_2^{(l)} = \frac{1}{x-\lambda_k} \cdot \underbrace{M_l \frac{x(x-1) \prod_{m \neq l} (x-t_m)}{\prod_{m \neq k} (x-\lambda_m)}}_{=: B_l}
\end{equation*}
より、
\begin{equation*}
    \lim_{x \to \lambda_k} B_l = M_l \frac{\lambda_k(\lambda_k-1) \prod_{m \neq l} (\lambda_k-t_m)}{\prod_{m \neq k} (\lambda_k-\lambda_m)} = M_l \frac{T(\lambda_k)}{L'(\lambda_k)(\lambda_k-t_l)} = M_l M^{k,l}
\end{equation*}
よって、
\begin{equation*}
    a_2^{(l)} = M_l \left[ \frac{M^{k,l}}{u} + \sum_{p=0}^\infty M^{k,l,p} u^p \right]
\end{equation*}
また、
\begin{align*}
    r(x,t) &= \frac{3}{4} u^{-2} - \nu_k u^{-1} + \sum_{p=0}^\infty r_{k,p} u^p \\
    \frac{\partial r}{\partial t_l} &= \frac{\partial}{\partial t_l} \left( \frac{3}{4} u^{-2} \right) - \frac{\partial}{\partial t_l} \left( \nu_k u^{-1} \right) + \sum_{p=0}^\infty \frac{\partial}{\partial t_l} \left( r_{k,p} u^p \right) \notag \\
    &= \frac{3}{2} \frac{\partial \lambda_k}{\partial t_l} u^{-3} - \nu_k \frac{\partial \lambda_k}{\partial t_l} u^{-2} - \frac{\partial \nu_k}{\partial t_l} u^{-1} + \sum_{p=0}^\infty \left( \frac{\partial r_{k,p}}{\partial t_l} - (p+1) r_{k,p+1} \frac{\partial \lambda_k}{\partial t_l} \right) u^p
\end{align*}
と展開すれば
\begin{align*}
    \frac{1}{2} a_2^{(l)} \frac{\partial^3 a_2^{(l)}}{\partial x^3} &= \frac{1}{2} \left( M_l M^{k,l} u^{-1} + \mathcal{O}(1) \right) \left( -6 M_l M^{k,l} u^{-4} + \mathcal{O}(1) \right) \\
    &= -3 M_l^2 \left(M^{k,l}\right)^2 u^{-5} + \mathcal{O}(u^{-4}) \\[8pt]
    \frac{\partial}{\partial x} \left( r(x,t) (a_2^{(l)})^2 \right) &= \frac{\partial}{\partial x} \left[ \left( \frac{3}{4} u^{-2} + \mathcal{O}(1) \right) \left( M_l M^{k,l} u^{-1} + \mathcal{O}(1) \right)^2 \right] \\
    &= \frac{\partial}{\partial x} \left( \frac{3}{4} M_l^2 \left(M^{k,l}\right)^2 u^{-4} + \mathcal{O}(u^{-3}) \right) \\
    &= -3 M_l^2 \left(M^{k,l}\right)^2 u^{-5} + \mathcal{O}(u^{-4}) \\[8pt]
    a_2^{(l)} \frac{\partial r}{\partial t_l} &\sim \mathcal{O}(u^{-1}) \mathcal{O}(u^{-3}) \sim \mathcal{O}(u^{-4})
\end{align*}

よって、
\begin{equation*}
    A_l \sim \mathcal{O}(u^{-4})
\end{equation*}

より，(2)も成り立つ．
\end{proof}

\vspace{1em}
以上のことから，有理関数 $A_l(x)$ は
\[
    A_l(x) = \frac{W_l}{x-t_l} + \sum_{k=1}^g \sum_{q=1}^4 \frac{W_q^{k,l}}{(x-\lambda_k)^q}
\]
と書かれる．ここで
\begin{equation} \tag{3.93}
    W_q^{k,l} = 0 \qquad (k, l = 1, \cdots, g), \quad (q = 1, 2, 3, 4)
\end{equation}
が成立するならば、補題3.11(4)から必然的に $W_l=0$ であり，(実際、$x=u^{-1}$を代入しても1位の零点しか持たない。)$A_l(x)=0$ が従う．そこで条件(3.93)を1つ1つ書き下していく．

\begin{align*}
    &W_4^{k,l} = 0 \qquad &&\frac{\partial \lambda_k}{\partial t_l} - M_l \left[ 2M^{k,l}\nu_k - M^{k,l,0} \right] = 0 \\[8pt]
    &W_3^{k,l} = 0 \qquad &&r_{k,0} - \nu_k^2 = 0 \\[8pt]
    &W_2^{k,l} = 0 \qquad &&\frac{\partial \nu_k}{\partial t_l} - M_l \left[ M^{k,l}r_{k,1} + M^{k,l,1}\nu_k - \frac{3}{2}M^{k,l,2} \right] = 0
\end{align*}

そして最後に $W_1^{k,l}=0$ から
\[
    M^{k,l} \frac{\partial}{\partial t_l} r_{k,0} - M^{k,l,0} \frac{\partial \nu_k}{\partial t_l} = \left( M^{k,l}r_{k,1} + M^{k,l,1}\nu_k - \frac{3}{2}M^{k,l,2} \right) \frac{\partial \lambda_k}{\partial t_l}
\]
を得る．
\end{story}

与えられた $A_l(x)$ を $u$ のべき級数として直接展開し、係数を整理する。
\begin{align*}
    A_l(x) &= \frac{1}{2} \left( M_l M^{k,l} u^{-1} + \sum_{p=0}^\infty M_l M^{k,l,p} u^p \right) \left( -6 M_l M^{k,l} u^{-4} + 6 M_l M^{k,l,3} + \mathcal{O}(u) \right) \\
    &\quad - \frac{\partial}{\partial x} \left\{ \left( \frac{3}{4} u^{-2} - \nu_k u^{-1} + \sum_{p=0}^\infty r_{k,p} u^p \right) \left( M_l^2 (M^{k,l})^2 u^{-2} + 2M_l^2 M^{k,l} M^{k,l,0} u^{-1} \right. \right. \\
    &\hspace{4em} \left. \left. + (M_l^2(M^{k,l,0})^2 + 2M_l^2 M^{k,l} M^{k,l,1}) + (2M_l^2 M^{k,l} M^{k,l,2} + 2M_l^2 M^{k,l,0} M^{k,l,1}) u + \mathcal{O}(u^2) \right) \right\} \\
    &\quad + \left( M_l M^{k,l} u^{-1} + \sum_{p=0}^\infty M_l M^{k,l,p} u^p \right) \left( \frac{3}{2} \frac{\partial \lambda_k}{\partial t_l} u^{-3} - \nu_k \frac{\partial \lambda_k}{\partial t_l} u^{-2} - \frac{\partial \nu_k}{\partial t_l} u^{-1} + \frac{\partial r_{k,0}}{\partial t_l} - r_{k,1}\frac{\partial \lambda_{k}}{\partial t_l} + \mathcal{O}(u) \right) \\
    &= -3 M_l^2 M^{k,l} M^{k,l,0} u^{-4} - 3 M_l^2 M^{k,l} M^{k,l,1} u^{-3} - 3 M_l^2 M^{k,l} M^{k,l,2} u^{-2} \\
    &\quad + \left( \underline{-3 M_l^2 M^{k,l} M^{k,l,3}} + \underline{3 M_l^2 M^{k,l} M^{k,l,3}} \right) u^{-1} \\
    &\quad - \left\{ -3 \left( \frac{3}{2} M_l^2 M^{k,l} M^{k,l,0} - \nu_k M_l^2 (M^{k,l})^2 \right) u^{-4} \right. \\
    &\hspace{3em} - 2 \left( \frac{3}{4} M_l^2 (M^{k,l,0})^2 + \frac{3}{2} M_l^2 M^{k,l} M^{k,l,1} - 2\nu_k M_l^2 M^{k,l} M^{k,l,0} + r_{k,0} M_l^2 (M^{k,l})^2 \right) u^{-3} \\
    &\hspace{3em} \left. - \left( \frac{3}{2} M_l^2 M^{k,l} M^{k,l,2} + \frac{3}{2} M_l^2 M^{k,l,0} M^{k,l,1} - \nu_k M_l^2 (M^{k,l,0})^2 - 2\nu_k M_l^2 M^{k,l} M^{k,l,1} \right. \right.\\
    &\hspace{3em}\left. \left.+ 2 r_{k,0} M_l^2 M^{k,l} M^{k,l,0} + r_{k,1} M_l^2 (M^{k,l})^2 \right) u^{-2} \right\} \\
    &\quad + \frac{3}{2} M_l M^{k,l} \frac{\partial \lambda_k}{\partial t_l} u^{-4} + \left( - M_l M^{k,l} \nu_k \frac{\partial \lambda_k}{\partial t_l} + \frac{3}{2} M_l M^{k,l,0} \frac{\partial \lambda_k}{\partial t_l} \right) u^{-3} \\
    &\quad + \left( - M_l M^{k,l} \frac{\partial \nu_k}{\partial t_l} - M_l M^{k,l,0} \nu_k \frac{\partial \lambda_k}{\partial t_l} + \frac{3}{2} M_l M^{k,l,1} \frac{\partial \lambda_k}{\partial t_l} \right) u^{-2} \\
    &\quad + \left( M_l M^{k,l} \frac{\partial r_{k,0}}{\partial t_l} - M_l M^{k,l} r_{k,1}\frac{\partial \lambda_{k}}{\partial t_l} - M_l M^{k,l,0} \frac{\partial \nu_k}{\partial t_l} - M_l M^{k,l,1} \nu_k \frac{\partial \lambda_k}{\partial t_l} + \frac{3}{2} M_l M^{k,l,2} \frac{\partial \lambda_k}{\partial t_l} \right) u^{-1} + \mathcal{O}(1)
\end{align*}
これを整理し、各 $u^{-j}$ の係数 $W_j^{k,l}$ を取り出して $=0$ と置く。

\paragraph{1. $u^{-4}$ の係数 $W_4^{k,l}$:}
\begin{align*}
    W_4^{k,l} &= -3 M_l^2 M^{k,l} M^{k,l,0} + \frac{9}{2} M_l^2 M^{k,l} M^{k,l,0} - 3 \nu_k M_l^2 (M^{k,l})^2 + \frac{3}{2} M_l M^{k,l} \frac{\partial \lambda_k}{\partial t_l} \\
    &= \frac{3}{2} M_l M^{k,l} \left( M_l \left[ M^{k,l,0} - 2\nu_k M^{k,l} \right] + \frac{\partial \lambda_k}{\partial t_l} \right) = 0.
\end{align*}
したがって、$\frac{\partial \lambda_k}{\partial t_l} = M_l (2 \nu_k M^{k,l} - M^{k,l,0})$ を得る。

\paragraph{2. $u^{-3}$ の係数 $W_3^{k,l}$:}
項のキャンセルを丁寧に整理すると、
\begin{align*}
    W_3^{k,l} &= -3 M_l^2 M^{k,l} M^{k,l,1} + \frac{3}{2} M_l^2 (M^{k,l,0})^2 + 3 M_l^2 M^{k,l} M^{k,l,1} - 4 \nu_k M_l^2 M^{k,l} M^{k,l,0} + 2r_{k,0} M_l^2 (M^{k,l})^2 \\
    &\quad - M_l M^{k,l} \nu_k \frac{\partial \lambda_k}{\partial t_l} + \frac{3}{2} M_l M^{k,l,0} \frac{\partial \lambda_k}{\partial t_l} \\
    &= M_l^2 \left( \frac{3}{2} (M^{k,l,0})^2 - 4\nu_k M^{k,l} M^{k,l,0} + 2r_{k,0}(M^{k,l})^2 \right) + M_l \left( \frac{3}{2} M^{k,l,0} - M^{k,l}\nu_k \right) \frac{\partial \lambda_k}{\partial t_l} \\
    &= 2 M_l^2 (M^{k,l})^2 \left( r_{k,0} - \nu_k^2 \right) = 0.
\end{align*}
したがって、$r_{k,0} = \nu_k^2$ を得る。

\paragraph{3. $u^{-2}$ の係数 $W_2^{k,l}$:}
各項を展開し、$r_{k,0} = \nu_k^2$ を用いて整理する。
\begin{align*}
W_2^{k,l} &= -3 M_l^2 M^{k,l} M^{k,l,2} + \frac{3}{2} M_l^2 M^{k,l} M^{k,l,2} + \frac{3}{2} M_l^2 M^{k,l,0} M^{k,l,1} - \nu_k M_l^2 (M^{k,l,0})^2\\
&\hspace{1em}- 2\nu_k M_l^2 M^{k,l} M^{k,l,1} + 2r_{k,0} M_l^2 M^{k,l} M^{k,l,0} + r_{k,1} M_l^2 (M^{k,l})^2\\
&\hspace{1em}- M_l M^{k,l} \frac{\partial \nu_k}{\partial t_l} - M_l M^{k,l,0} \nu_k \frac{\partial \lambda_k}{\partial t_l} + \frac{3}{2} M_l M^{k,l,1} \frac{\partial \lambda_k}{\partial t_l}
\end{align*}
ここで $\frac{\partial \lambda_k}{\partial t_l}$ を含む項をまとめ、残りの $M_l^2$ の項と分ける。
\begin{align*}
&= M_l^2 \left( \frac{3}{2} M^{k,l,1} - M^{k,l,0} \nu_k \right) \left( \frac{1}{M_l} \frac{\partial \lambda_k}{\partial t_l} \right)\\
&+ M_l^2 \left( -\frac{3}{2} M^{k,l} M^{k,l,2} + \frac{3}{2} M^{k,l,0} M^{k,l,1} - \nu_k (M^{k,l,0})^2 - 2\nu_k M^{k,l} M^{k,l,1} + 2\nu_k^2 M^{k,l} M^{k,l,0} \right)\\
&+ r_{k,1} M_l^2 (M^{k,l})^2 - M_l M^{k,l} \frac{\partial \nu_k}{\partial t_l}
\end{align*}
第1項の括弧に $\frac{1}{M_l} \frac{\partial \lambda_k}{\partial t_l} = 2\nu_k M^{k,l} - M^{k,l,0}$ を代入して展開する。
\begin{align*}
&\left( \frac{3}{2} M^{k,l,1} - M^{k,l,0} \nu_k \right) ( 2\nu_k M^{k,l} - M^{k,l,0} )\\
&= 3\nu_k M^{k,l} M^{k,l,1} - \frac{3}{2} M^{k,l,0} M^{k,l,1} - 2\nu_k^2 M^{k,l} M^{k,l,0} + \nu_k (M^{k,l,0})^2
\end{align*}
これを元の式に代入し、相殺する項を整理する。
\begin{align*}
W_2^{k,l} &= M_l^2 \left( 3\nu_k M^{k,l} M^{k,l,1} - \frac{3}{2} M^{k,l,0} M^{k,l,1} - 2\nu_k^2 M^{k,l} M^{k,l,0} + \nu_k (M^{k,l,0})^2 \right.\\
&\left. -\frac{3}{2} M^{k,l} M^{k,l,2} + \frac{3}{2} M^{k,l,0} M^{k,l,1} - \nu_k (M^{k,l,0})^2 - 2\nu_k M^{k,l} M^{k,l,1} + 2\nu_k^2 M^{k,l} M^{k,l,0} \right)\\
&+ r_{k,1} M_l^2 (M^{k,l})^2 - M_l M^{k,l} \frac{\partial \nu_k}{\partial t_l}
\end{align*}

括弧内の相殺（$-\frac{3}{2}$ と $+\frac{3}{2}$、$-2\nu_k^2$ と $+2\nu_k^2$、$+\nu_k$ と $-\nu_k$）を実行し、残った $3\nu_k M^{k,l} M^{k,l,1}$ と $-2\nu_k M^{k,l} M^{k,l,1}$ を足し合わせる。
\begin{equation}
= M_l^2 \left( \nu_k M^{k,l} M^{k,l,1} - \frac{3}{2} M^{k,l} M^{k,l,2} \right) + r_{k,1} M_l^2 (M^{k,l})^2 - M_l M^{k,l} \frac{\partial \nu_k}{\partial t_l}
\end{equation}

全体を $-M_l M^{k,l}$ でくくる。

\begin{equation}
= - M_l M^{k,l} \left( \frac{\partial \nu_k}{\partial t_l} - M_l \left( \nu_k M^{k,l,1} - \frac{3}{2} M^{k,l,2} + r_{k,1} M^{k,l} \right) \right) = 0
\end{equation}
したがって、次式を得る。
\begin{equation}
\frac{\partial \nu_k}{\partial t_l} = M_l \left( \nu_k M^{k,l,1} - \frac{3}{2} M^{k,l,2} + r_{k,1} M^{k,l} \right)
\end{equation}

\paragraph{4. $u^{-1}$ の係数 $W_1^{k,l}$:}
\begin{align*}
    W_1^{k,l} &= M_l \left( M^{k,l} \frac{\partial r_{k,0}}{\partial t_l} + M^{k,l} r_{k,1} - M^{k,l,0} \frac{\partial \nu_k}{\partial t_l} - M^{k,l,1} \nu_k \frac{\partial \lambda_k}{\partial t_l} + \frac{3}{2} M^{k,l,2} \frac{\partial \lambda_k}{\partial t_l} \right) = 0.
\end{align*}
つまり、
\begin{equation*}
M^{k,l} \frac{\partial r_{k,0}}{\partial t_l} - M^{k,l,0} \frac{\partial \nu_k}{\partial t_l} = \left(M^{k,l} r_{k,1} + M^{k,l,1} \nu_k  - \frac{3}{2} M^{k,l,2}\right) \frac{\partial \lambda_k}{\partial t_l}
\end{equation*}
以上より、各係数の消失からシステムを記述する関係式が導出された。

\begin{story}
まず，$W_3^{k,l}=0$ は，2階線型常微分方程式(3.25)の特異点 $x=\lambda_k$ が非対数的である，という条件と同一であり，新しい束縛条件は与えない．
\end{story}
$\phi(x) = u^s \sum_{n=0}^\infty a_n u^n$ として、微分方程式 $\frac{d^2\phi(x)}{dx^2} = r(x,t)\phi(x)$ に代入する。
\begin{equation*}
    u^s \sum_{n=0}^\infty (s+n)(s+n-1) a_n u^{n-2} = \left( \frac{3}{4} u^{-2} - \nu_k u^{-1} + \sum_{n=0}^\infty r_{k,n} u^n \right) \left( u^s \sum_{n=0}^\infty a_n u^n \right)
\end{equation*}
両辺の $u$ の各べきの係数を比較する。

\paragraph{1. $u^{s-2}$ の係数:}
\begin{align*}
    s(s-1) a_0 &= \frac{3}{4} a_0 \\
    \therefore \quad \left(s + \frac{1}{2}\right)\left(s - \frac{3}{2}\right) a_0 &= 0
\end{align*}
$a_0 \neq 0$ より $s = -1/2, 3/2$ を得る。対数項の有無を調べるため、小さい根 $s = -1/2$ を採用して高次の係数を比較する。

\paragraph{2. $u^{-3/2}$ の係数（$s=-1/2, n=1$ の項）:}
\begin{align*}
    \left(\frac{1}{2}\right)\left(-\frac{1}{2}\right) a_1 &= \frac{3}{4} a_1 - \nu_k a_0 \\
    -\frac{1}{4} a_1 &= \frac{3}{4} a_1 - \nu_k a_0 \\
    \therefore \quad a_1 &= \nu_k a_0
\end{align*}

\paragraph{3. $u^{-1/2}$ の係数（$s=-1/2, n=2$ の項）:}
\begin{align*}
    \left(\frac{3}{2}\right)\left(\frac{1}{2}\right) a_2 &= \frac{3}{4} a_2 - \nu_k a_1 + r_{k,0} a_0 \\
    \frac{3}{4} a_2 &= \frac{3}{4} a_2 - \nu_k a_1 + r_{k,0} a_0 \\
    \therefore \quad 0 &= -\nu_k a_1 + r_{k,0} a_0
\end{align*}
先ほど求めた $a_1 = \nu_k a_0$ を代入して整理すると、
\begin{equation*}
    (\nu_k^2 - r_{k,0}) a_0 = 0
\end{equation*}
$a_0 \neq 0$ より、$r_{k,0} = \nu_k^2$ が導かれる。
\begin{story}
また $W_1^{k,l}=0$ は，$W_4^{k,l}=W_2^{k,l}=0$ から $W_1^{k,l}$の係数部分に注目して、代入すると、
\[
    M^{k,l} \frac{\partial}{\partial t_l} r_{k,0} - \left( 2M^{k,l}\nu_k - \frac{1}{M_l} \frac{\partial \lambda_k}{\partial t_l} \right) \frac{\partial \nu_k}{\partial t_l} = \frac{1}{M_l} \frac{\partial \nu_k}{\partial t_l} \cdot \frac{\partial \lambda_k}{\partial t_l}
\]
すなわち次のようにも表され，これも非対数的条件 $W_3^{k,l}=0$ に含まれている．($W_3^{k,l}=0$を$t_{l}$で微分すれば良い。)
\begin{equation}
    \frac{\partial}{\partial t_l} r_{k,0} - 2\nu_k \frac{\partial \nu_k}{\partial t_l} = 0 \label{eq:rk0_derivative}
\end{equation}

したがって、本質的な条件は方程式系 $W_4^{k,l} = W_2^{k,l} = 0$ である。ここで、ハミルトニアン $K_l$ の具体形(3.86)を用いて直接計算で確認すれば、この方程式系がハミルトン系(3.91)に他ならないことがわかる。この部分の確認は読者にお任せする。以上の説明により補題3.9が確かに成り立つことを認め、先に進むことにしよう。
\end{story}

実際、
運動量の変数変換式において、括弧内を $\omega_k$ と定義する。
\begin{equation*}
    \nu_k = \mu_k + \frac{1}{2} \underbrace{\left( \frac{1-\kappa_0}{\lambda_k} + \frac{1-\kappa_1}{\lambda_k-1} + \sum_{i=1}^g \frac{1-\theta_i}{\lambda_k-t_i} - \sum_{m \neq k} \frac{1}{\lambda_k-\lambda_m} \right)}_{=: \omega_k}
\end{equation*}
元のハミルトニアン $H_l$ とパラメータ $A_{k,l}$ は次で与えられる。
\begin{align*}
    H_l &= M_l \left[ \sum_{k=1}^g M^{k,l} \left\{ \mu_k^2 - A_{k,l} \mu_k + \frac{\kappa}{\lambda_k(\lambda_k-1)} \right\} \right] \\
    A_{k,l} &= \frac{\kappa_0}{\lambda_k} + \frac{\kappa_1}{\lambda_k-1} + \sum_{i=1}^g \frac{\theta_i - \delta_{il}}{\lambda_k-t_i}
\end{align*}
$\mu_k = \nu_k - \frac{1}{2}\omega_k$ を $H_l$ に代入して、
\begin{equation*}
    H_l = M_l \left[ \sum_{k=1}^g M^{k,l} \left\{ \left(\nu_k - \frac{1}{2}\omega_k\right)^2 - A_{k,l} \left(\nu_k - \frac{1}{2}\omega_k\right) + \frac{\kappa}{\lambda_k(\lambda_k-1)} \right\} \right]
\end{equation*}
ここで、$K_l$ は $H_l$ に $\nu_k$ を含まない項を足して定義されている。
\begin{equation*}
    K_l = H_l + \frac{1-\theta_l}{2} \left( \frac{1-\kappa_0}{t_l} + \frac{1-\kappa_1}{t_l-1} + \sum_{m \neq l} \frac{1-\theta_m}{t_l-t_m} - \sum_{k=1}^g \frac{1}{t_l-\lambda_k} \right)
\end{equation*}
したがって、$\nu_k$ による偏微分は以下のようになる。
\begin{align*}
    \frac{\partial K_l}{\partial \nu_k} &= \frac{\partial H_l}{\partial \nu_k} = M_l M^{k,l} \left\{ 2\left(\nu_k - \frac{1}{2}\omega_k\right) - A_{k,l} \right\} \\
    &= M_l \left( 2 M^{k,l} \nu_k - M^{k,l}(\omega_k + A_{k,l}) \right)
\end{align*}
これまでの特異点解析より $\frac{\partial \lambda_k}{\partial t_l} = M_l(2 M^{k,l} \nu_k - M^{k,l,0})$ であるため、
\begin{equation*}
    M^{k,l,0} = M^{k,l}(\omega_k + A_{k,l})
\end{equation*}
を示せばよい。$\omega_k$ と $A_{k,l}$ を足し合わせると、パラメータ $\kappa, \theta$ が相殺し、
\begin{equation*}
    \omega_k + A_{k,l} = \frac{1}{\lambda_k} + \frac{1}{\lambda_k-1} + \sum_{h \neq l} \frac{1}{\lambda_k-t_h} - \sum_{m \neq k} \frac{1}{\lambda_k-\lambda_m}
\end{equation*}
となる。一方で、ローラン展開の定義 $a_2^{(l)}(x,t) = M_l(M^{k,l}u^{-1} + M^{k,l,0} + \mathcal{O}(u))$ より、
\begin{align*}
    a_2^{(l)}(x,t) - M_l M^{k,l} u^{-1} &= M_l \left( \frac{T(x)}{L(x)(x-t_l)} - \frac{T(\lambda_k)}{(x-\lambda_k)L'(\lambda_k)(\lambda_k-t_l)} \right) \\
    &= \frac{M_l}{u} \left( \frac{x(x-1)\prod_{i \neq l}(x-t_i)}{\prod_{h \neq k}(x-\lambda_h)} - \frac{\lambda_k(\lambda_k-1)\prod_{i \neq l}(\lambda_k-t_i)}{\prod_{h \neq k}(\lambda_k-\lambda_h)} \right)
\end{align*}
ここで、関数 $F(x) = \frac{x(x-1)\prod_{i \neq l}(x-t_i)}{\prod_{h \neq k}(x-\lambda_h)}$ とおくと、
\begin{equation*}
    M_l M^{k,l,0} = \lim_{x \to \lambda_k} M_l \frac{F(x) - F(\lambda_k)}{x-\lambda_k} = M_l F'(\lambda_k)
\end{equation*}
となるため、$M^{k,l,0} = F'(\lambda_k)$ である。対数微分法を用いると、
\begin{align*}
    \frac{F'(x)}{F(x)} &= \frac{1}{x} + \frac{1}{x-1} + \sum_{i \neq l} \frac{1}{x-t_i} - \sum_{h \neq k} \frac{1}{x-\lambda_h} \\
    \therefore \quad F'(\lambda_k) &= \left( \frac{1}{\lambda_k} + \frac{1}{\lambda_k-1} + \sum_{i \neq l} \frac{1}{\lambda_k-t_i} - \sum_{h \neq k} \frac{1}{\lambda_k-\lambda_h} \right) F(\lambda_k) \\
    &= M^{k,l}(\omega_k + A_{k,l})
\end{align*}
以上より、$M^{k,l,0} = M^{k,l}(\omega_k + A_{k,l})$ が示された。
よって、
\begin{equation*}
    \frac{\partial \lambda_k}{\partial t_l} = \frac{\partial K_l}{\partial \nu_k}
\end{equation*}
が成立する。($\frac{\partial \nu_k}{\partial t_l} = \frac{\partial K_l}{\partial \lambda_k}$は面倒なのでパス)